%----------------------------------------------------------------------------------------
%	PACKAGES AND OTHER DOCUMENT CONFIGURATIONS
%----------------------------------------------------------------------------------------

\documentclass[
	11pt, % Default font size, select one of 10pt, 11pt or 12pt
	fleqn, % Left align equations
	a4paper, % Paper size, use either 'a4paper' for A4 size or 'letterpaper' for US letter size
	%oneside, % Uncomment for oneside mode, this doesn't start new chapters and parts on odd pages (adding an empty page if required), this mode is more suitable if the book is to be read on a screen instead of printed
]{LegrandOrangeBook}

% Book information for PDF metadata, remove/comment this block if not required 
\hypersetup{
	pdftitle={Fisica - Meccanica}, % Title field
	pdfauthor={DM}, % Author field
	pdfsubject={Meccanica}, % Subject field
	pdfkeywords={Keyword1, Keyword2, ...}, % Keywords
	pdfcreator={LaTeX}, % Content creator field
}

\addbibresource{sample.bib} % Bibliography file

\definecolor{ocre}{RGB}{243, 102, 25} % Define the color used for highlighting throughout the book

\chapterimage{orange1.jpg} % Chapter heading image
\chapterspaceabove{6.5cm} % Default whitespace from the top of the page to the chapter title on chapter pages
\chapterspacebelow{6.75cm} % Default amount of vertical whitespace from the top margin to the start of the text on chapter pages

%----------------------------------------------------------------------------------------

\begin{document}

%----------------------------------------------------------------------------------------
%	TITLE PAGE
%----------------------------------------------------------------------------------------

\titlepage % Output the title page
	{\includegraphics[width=\paperwidth]{background.pdf}} % Code to output the background image, which should be the same dimensions as the paper to fill the page entirely; leave empty for no background image
	{ % Title(s) and author(s)
                \centering
		\vspace{16pt} % Vertical whitespace
		{\LARGE Scuole superiori\par} % Subtitle
		\sffamily % Font styling
		{\Huge\bfseries Meccanica\par} % Book title
		\vspace{24pt} % Vertical whitespace
		{\huge\bfseries OSB\par} % Author name
	}

%----------------------------------------------------------------------------------------
%	COPYRIGHT PAGE
%----------------------------------------------------------------------------------------

\thispagestyle{empty} % Suppress headers and footers on this page

~\vfill % Push the text down to the bottom of the page

\noindent Copyright \copyright\ 2023 OSB\\ % Copyright notice

\noindent \textsc{Published by OSB}\\ % Publisher

% \noindent \textsc{\href{https://www.latextemplates.com/template/legrand-orange-book}{book-website.com}}\\ % URL

\noindent Licensed under the Creative Commons Attribution-NonCommercial 4.0 License (the ``License''). You may not use this file except in compliance with the License. You may obtain a copy of the License at \url{https://creativecommons.org/licenses/by-nc-sa/4.0}. Unless required by applicable law or agreed to in writing, software distributed under the License is distributed on an \textsc{``as is'' basis, without warranties or conditions of any kind}, either express or implied. See the License for the specific language governing permissions and limitations under the License.\\ % License information, replace this with your own license (if any)

\noindent \textit{Latest version \today} % Printing/edition date

%----------------------------------------------------------------------------------------
%	TABLE OF CONTENTS
%----------------------------------------------------------------------------------------

\pagestyle{empty} % Disable headers and footers for the following pages

\tableofcontents % Output the table of contents

\listoffigures % Output the list of figures, comment or remove this command if not required

\listoftables % Output the list of tables, comment or remove this command if not required

\pagestyle{fancy} % Enable default headers and footers again

\cleardoublepage % Start the following content on a new page


----------------------------------------------------------------------------------------
%	PART
%----------------------------------------------------------------------------------------

\part{Meccanica}

\part{Storia}

\section{Cronologia}
\paragraph{\dots}
\paragraph{XVI secolo}
\begin{itemize}
  \item Nepier (Nepero) introduce i logaritmi
\end{itemize}
\paragraph{XVII secolo}
\begin{itemize}
  \item Fermat
  \item Descartes (Cartesio) illustra ne \textit{La Gèometrie} i fondamenti della \textbf{geometria analitica}
  \item Huygens, Pascal
  \item Newton e Leibniz sviluppano contemporaneamente i fondamenti del \textbf{calcolo differenziale} e \textbf{integrale}, nell'ambito dello studio della \textbf{dinamica}
  \item fratelli Johann e Jakob Bernoulli
  \item de l'Hopital, Taylor
\end{itemize}
\paragraph{XVIII secolo}
\begin{itemize}
  \item Euler (Eulero): analisi matematica; soluzione equazioni differenziali; teoria dei numeri; analisi complessa (estensione di funzioni reali in ambito complesso; identità di Eulero); topologia e teoria dei grafi (problema die 7 ponti di Konigsberg)
  \item d'Alembert si dedica allo studio del moto dei corpi e lla meccanica razionale
  \item Legendre
  \item Bayes: probabilità
  \item istituzione di scuole scientifiche, Parigi importante centro scientifico del tempo
  \item Laplace: meccanica razionale e celeste (\textit{Mécanique Céleste}); trasformata; calcolo differenziale: potenziale, laplaciano ed equazione di Laplace
  \item Lagrange: formulazione lagrangiana della meccanica (\textit{Mécanique analytique}); calcolo delle variazioni; metodo dei moltiplicatori di Lagrange; teoria dei numeri
\end{itemize}
\paragraph{XIX secolo}
\begin{itemize}
  \item Jacobi: algebra lineare (determinante di matrici)
  \item Cauchy: algebra lineare; analisi complessa; statistica; teoria dei numeri; meccanica dei solidi
  \item Fourier: studio della trasmissione del calore; serie e trasformata di Fourier
  \item Gauss: teorema fondamentale dell'algebra; teoria dei numeri
  \item Dirichlet:
  \item Riemann: teoria dei numeri; geometria
  \item Hamilton: quaternioni; algebra lineare (teorema di Cayley-Hamilton); riformulazione della meccanica lagrangiana nella meccanica hamiltoniana
  \item Wierestrass: definizione rigorosa dei fondamenti dell'analisi (teorema di Weierstrass su esisteanza di minimi e massimi di funzioni a variabile reale)
  \item Boole: algebra sugli insiemi, logica, e teoria dell'informazione
  \item Peano: tentativo di definizione assiomatica della matematica
  \item Cantor: studio degli insiemi infiniti e la loro dimensione
\end{itemize}
\paragraph{XX secolo}
\begin{itemize}
  \item la matematica della probabilità e della meccanica quantistica: Lebesgue, Hilbert, von Neumann, Kolmogorov
  \item la nascita dell'informatica: Turing, Von Neumann
  \item la teoria dell'informazione: Shannon
  \item la teoria dei giochi: von Neumann, Morgestern e Nash
  \item l'incompletezza della matematica: Godel
\end{itemize}



\part{Insiemistica e Logica}

\part{Algebra in $\mathbb{R}$}

\chapter{Algebra simbolica - Calcolo letterale}

% ------------------------------------------------------------------------------
\section{Monomi}
\begin{definition}[Monomio] Un monomio è un'espressione matematica costituita dal prodotto di un coefficiente esplicitamente numerico e una parte letterale, nella quale compaiono unicamente moltiplicazioni e potenze intere.
\end{definition}
\begin{remark} Viene richiesto che le potenze della parte letterale siano intere, per evitare di porre delle condizioni sulle basi delle potenze, essendo le potenze non intere definite solo per numeri reali positivi.
\end{remark}

\begin{definition}[Monomi simili] I monomi simili sono i monomi che hanno la stessa parte letterale.
\end{definition}

\subsection{Somma e differenza}
\subsection{Prodotto e divisione}
\subsection{Potenze e radici}
\begin{definition}[Potenze e radici intere] La potenza intera di ordine $n \in \mathbb{N}$ di un monomio $x$ è definita come il prodotto di $x$ per se stesso $n$ volte,
\begin{equation}
  p_n(x) := x^n = \underbrace{x \cdot x \cdot \dots \cdot x}_{\text{$n$ volte}} \ .
\end{equation}
L'operazione inversa, quando possibile, è definita come radice di ordine $n$,
\begin{equation}
  x := \sqrt[n]{p_n(x)} = p_n(x)^{\frac{1}{n}} \ .
\end{equation}
\end{definition}
\begin{remark} Per una comprensione più completa, bisogna rifarsi all'algebra dei numeri complessi \ref{book:complex_algebra}.
\end{remark}

\begin{definition}[Potenze e radici non intere]
\end{definition}

% ------------------------------------------------------------------------------
\section{Polinomi}
\begin{definition}[Polinomio]
Un polinomio reale di grado $n$ viene definito come una combinazione lineare dei monomi di grado $\le n$,
\begin{equation}
 p(x) = a_0 + a_1 x + \dots + a_n x^n = \sum_{i=0}^{n} a_i x^i
\end{equation}
\end{definition}

\begin{definition}[Zero di un polinomio] Viene definito zero -- o radice -- (reale) di un polinomio un numero $\overline{x}$ ($\in \mathbb{R}$) tale che $p(\overline{x}) = 0$.
\end{definition}

\begin{definition}[Scomposizione in fattori]
    Ogni polinomio di variabile reale e coefficienti reali può essere scomposto in fattori polinomiali di primo e secondo grado. Ad esempio un polinomio di grado $n$, può essere scomposto nel prodotto di $p_1$ polinomi di primo grado e $p_2$ polinomi di secondo grado,
    \begin{equation}
        \begin{aligned}
            P(x) & = a_0 + a_1 x + \dots + a_n x^n = \\
            & = (A_1 x+B_1) \dots (A_{p_1}x+B_{p_1}) (C_1 x^2+D_1 x+E_1) \dots (C_{p_2} x^2 + D_{p_2} x + E_{p_2} ) = \\
            & = K (x+\tilde{B}_1) \dots (x+\tilde{B}_{p_1}) (x^2 + \tilde{D}_1 x + \tilde{E}_1) \dots (x^2 + \tilde{D}_{p_2} x + \tilde{E}_{p_2})
        \end{aligned}
    \end{equation}
    con $p_1 + 2 p_2 = n$.
\end{definition}

\subsection{Scomposizioni notevoli}
La dimostrazione di queste identità viene ricavata facilmente tramite calcolo diretto.
\paragraph{Quadrato della somma}
\begin{equation}
    (a+b)^2 = a^2 + 2 ab + b^2
\end{equation}
\paragraph{Cubo della somma}
\begin{equation}
    (a+b)^3 = a^3 + 3 a^2 b + 3 a b^2 + b^3
\end{equation}
\paragraph{Potenza di un binomio}
\begin{equation}
    (a+b)^n = \sum_{k=0}^{n} \binom{n}{k} a^{n-k} b^{k}
\end{equation}
\paragraph{Differenza di quadrati}
\begin{equation}
    a^2 - b^2 = (a+b)(a-b)
\end{equation}
\paragraph{Somma e differenza di cubi}
\begin{equation}
    a^3 \mp b^3 = (a \mp b)(a^2 \pm ab + b^2)
\end{equation}
\paragraph{Differenza di potenze di grado qualsiasi}
\begin{equation}
\begin{aligned}
    a^n - b^n & =
        (a-b) \sum_{k=0}^{n-1} a^{n-1-k} b^k = \\
    & = (a-b)(a^{n-1} + a^{n-2}b + a^{n-3}b^2 + \dots + a b^{n-2} + b^{n-1})
\end{aligned}
\end{equation}
\paragraph{Somma di potenze di grado dispari}
Per ogni numero dispari $n = 2m + 1$, $m \in \mathbb{N}$, si può dimostrare che
\begin{equation}
\begin{aligned}
    a^{2m+1} + b^{2m+1} & =
        (a + b) \sum_{k=0}^{2m} (-1)^k a^{n-1-k} b^k = \\
    & = (a + b)(a^{n-1} - a^{n-2}b + a^{n-3}b^2 + \dots - a b^{n-2} + b^{n-1})
\end{aligned}
\end{equation}

\subsection{Divisione tra polinomi e resto}
\begin{definition}[Divisione tra polinomi]
    Dati due polinomi $P(x)$, $D(x)$ esistono due polinomi $Q(x)$, $R(x)$ tali che
    \begin{equation}
        P(x) = D(x) Q(x) + R(x) \ ,
    \end{equation}
    con il grado di $R(x)$ minore del grado di $Q(x)$. I polinomi $Q(x)$ e $R(x)$ vengono definiti rispettivamente \textbf{quoziente} e \textbf{resto} della divisione.
\end{definition}

\subsubsection{Divisione esatta e scomposizone di polinomi: regola di Ruffini}

\subsection{Formule di Viete e Newton}
{\color{red} Relazioni tra i coefficienti di un polinomio e le sue radici}


% ------------------------------------------------------------------------------
\section{Potenze e radici}
\begin{definition}[Potenze e radici intere]
\end{definition}

% ------------------------------------------------------------------------------
\section{Esponenziali e logaritmi}
\subsection{Esponenziale}
\begin{definition}[Esponenziale] L'elevamento a potenza di un numero $a$
\begin{equation}
   y = a^x
\end{equation}
è un operazione che coinvolge due numeri, $a$ detto base e $x$ detto esponente.
\end{definition}
\subsubsection{Potenze non intere, valori ammissibili}
\subsubsection{Proprietà}
\paragraph{Prodotto di potenze con la stessa base}
\begin{equation} a^m a^n = a^{m+n} \end{equation}
\paragraph{Potenza di potenza}
\begin{equation} (a^m)^n = a^{mn} \end{equation}
\paragraph{Prodotto di potenze con lo stesso esponente}
\begin{equation} a^m b^m = (ab)^{m} \end{equation}

% ------------------------------------------------------------------------------
\subsection{Logaritmo}
\begin{definition}[Logaritmo] Il logaritmo è l'operazione inversa
\begin{equation}
    x = \log_a y \qquad \text{se $y = a^x$}
\end{equation}
\end{definition}
\subsubsection{Potenze non intere, valori ammissibili}
\subsubsection{Proprietà}
\paragraph{Somma di logaritmi con la stessa base}
\begin{equation}
  \log_a m + \log_a n = \log_a (mn)
\end{equation}
Dimostrazione
\begin{equation}
\begin{cases}
  m = a^{\log_a m} \\
  n = a^{\log_a n} \\
  mn = a^{\log_a mn}
\end{cases}
\rightarrow
\begin{aligned}
    mn & = m \cdot n \\
    a^{\log_a mn} & = a^{\log_a m} a^{\log_a n} = a^{\log_a m + \log_a n}
\end{aligned}
\end{equation}
\begin{equation}
  \rightarrow \log_a mn = \log_a m + \log_a n
\end{equation}
\paragraph{Prodotto di un logaritmo per uno scalare}
\begin{equation}
  b \log_a m = \log_a m^{b}
\end{equation}
\paragraph{Cambio di base di un logaritmo}
\begin{equation}
     \log_b m = \log_b a \log_a m 
\end{equation}
Dimostrazione
\begin{equation}
\begin{cases}
    m = b^{\log_b m} \\
    a = b^{\log_b a} \\
    m = a^{\log_a m} = \left( b^{\log_b a} \right)^{\log_a m} = b^{\log_b a \, \log_a m}
\end{cases}
\end{equation}
e confrontando le due espressioni per $m$ si ottiene
\begin{equation}
    \rightarrow \log_b m = \log_b a \, \log_a m
\end{equation}

% ------------------------------------------------------------------------------
\section{Funzioni armoniche}
\subsection{La circonferenza e la definizione delle funzioni seno e coseno}

\subsection{La definizione delle funzioni tangente, cotangente, secante e cosecante}

\subsection{Formule del seno e coseno di somme e differenze}
\begin{equation}
\begin{aligned}
    \cos(\alpha \pm \beta) & = \cos(\alpha) \cos(\beta) \mp \sin(\alpha) \sin(\beta) \\
    \sin(\alpha \pm \beta) & = \sin(\alpha) \cos(\beta) \pm \cos(\alpha) \sin(\beta) \\
\end{aligned}
\end{equation}
\begin{equation}
\begin{aligned}
    \cos(\alpha) \cos(\beta) & = \dfrac{1}{2} \left[ \cos(\alpha + \beta) + \cos(\alpha-\beta) \right] \\
    \sin(\alpha) \sin(\beta) & = \dfrac{1}{2} \left[ \cos(\alpha + \beta) - \cos(\alpha-\beta) \right] \\
    \sin(\alpha) \cos(\beta) & = \dfrac{1}{2} \left[ \sin(\alpha + \beta) - \sin(\alpha-\beta) \right] \\
\end{aligned}
\end{equation}

% ------------------------------------------------------------------------------
\section{Funzioni iperboliche}


% ==============================================================================
\chapter{Equazioni}
\section{Equazioni algebriche}
\begin{definition}[Equazioni algebriche] \dots
\end{definition}

\subsection{Equazioni polinomiali}
\begin{definition}[Equazione polinomiale] Un'equazione polinomiale ha la forma
    \begin{equation} p(x) = 0 \end{equation}
dove $p(x)$ è un polinomio.
Il \textbf{grado} dell'equazione corrisponde al grado del polinomio $p(x)$, cioé alla potenza massima dei monomi. In maniera più esplicita, quindi, si può scrivere un'equazione polinomiale di grado $n$ come
    \begin{equation}
        a_n x^n + a_{n-1} x^{n-1} + \dots + a_1 x + a_0 = 0 \ , \qquad \text{con $a_n \ne 0$}
    \end{equation}
\end{definition}

\paragraph{Esistenza e numero delle soluzioni.} Un'equazione polinomiale di grado $n$ ha \textbf{al massimo} $n$ soluzioni reali. L'esistenza di soluzioni reali non è in generale garantita, mentre il \textbf{teorema fondamentale dell'algebra} assicura che esistano esattamente $n$ soluzioni complesse di un'equazione polinomiale con coefficienti complessi.

\subsubsection{Equazioni di primo grado}
La forma generale delle equazioni di primo grado è
\begin{equation}
    a x + b = 0 \ , \qquad \text{con $a \ne 0$}
\end{equation}
e la soluzione è
\begin{equation}
  x = -\dfrac{a_0}{a_1} \ .
\end{equation}
\subsubsection{Equazioni di secondo grado}
La forma generale delle equazioni di secondo grado è
\begin{equation}
    a x^2 + b x + c = 0 \ , \qquad \text{con $a \ne 0$}
\end{equation}
%
Un'equazione di secondo grado può ammettere nel campo dei numeri reali 2 soluzioni (distinte o coincidenti) o nessuna soluzione, a seconda del valore dell'espressione definita come \textbf{discriminante}, $\Delta := b^2 - 4 a c$:
\begin{itemize}
    \item $\Delta > 0$: due soluzioni reali distinte
    \item $\Delta = 0$: due soluzioni reali coincidenti
    \item $\Delta < 0$: nessuna soluzione reale
\end{itemize}
Quando il discriminante è non negativo, le soluzioni dell'equazione sono date dall'espressione
\begin{equation}
    x_{1,2} = \dfrac{-b \mp \sqrt{\Delta}}{2a} \ .
\end{equation}
%
\paragraph{Formula risolutiva dell'equazione di secondo grado.}
Una dimostrazione della formula risolutiva viene ricavata con la regola di completamento del quadrato
\begin{equation}
    \begin{aligned}
        0 & = a x^2 + b x + c = \\
          & = a x^2 + b x + \dfrac{b^2}{4a} - \dfrac{b^2}{4a} + c = \\
          & = a \left(x + \dfrac{b}{2a} \right)^2 - \dfrac{b^2}{4a} + c \\
    \end{aligned}
\end{equation}
\begin{equation}
    \rightarrow \left( x + \dfrac{b}{2a} \right)^2 = \dfrac{b^2}{4 a^2} - \dfrac{c}{a} = \dfrac{b^2 - 4 a c}{4a^2} = \dfrac{\Delta}{4a^2}
\end{equation}
\'E ora facile notare come questa equazione ha soluzioni solo quando il discriminante è non negativo. Quando il discriminante è non negativo, è possibile estrarre la radice quadra dell'espressione
\begin{equation}
    x_{1,2} + \dfrac{b}{2a} = \mp \dfrac{\sqrt{\Delta}}{2a} \qquad \rightarrow \qquad 
    x_{1,2} = \dfrac{-b  \mp \sqrt{\Delta}}{2a} \ . 
\end{equation}

\subsection{Equazioni algebriche razionali}
\begin{definition}[Equazioni algebriche razionali] Le equazioni algebriche razionali sono equazioni che contengono polinomi, loro rapporti e potenze intere.
\end{definition}
\begin{example}[Esempi di equazioni algebriche razionali]
\begin{equation}
    \dfrac{(x + 3)^2}{(x-1)} = 4 x 
\end{equation}
\begin{equation}
    \dfrac{2x}{x^2+1} = -\dfrac{1}{x}
\end{equation}
\end{example}

\subsubsection{Metodo di soluzione}
\begin{enumerate}
    \item Per prima cosa è necessario determinare le \textbf{condizioni di esistenza} di una soluzione. Poiché nelle equazioni può comparire la \textbf{divisione} tra polinomi, bisogna richiedere che questa e tutte le operazioni scritte nel problema abbiano senso: ad esempio, nelle condizioni di esistenza bisogna richiedere che non avvengano divisioni per zero.
    \item Successivamente, è possibile procedere con le semplificazioni per la ricerca della soluzione.
\end{enumerate}

\paragraph{Esempi}
Seguendo questo metodo di soluzione, procediamo a risolvere le equazioni dell'esempio {\color{red} cit}.
\begin{equation}
    \dfrac{(x + 3)^2}{(x-1)} = 4 x  \qquad \qquad \text{C.E.: } x \ne 1 \quad \rightarrow \quad x \in \mathbb{R}\backslash \{1\}
\end{equation}
\begin{equation}
    \begin{aligned}
        & (x + 3)^2 = 4x(x-1) \\
        & x^2+6x+9 = 4x^2 -4x \\
        & 3x^2-10x-9 = 0 \\
        & \rightarrow \quad x_{1,2} = \dfrac{5 \mp \sqrt{25+3\cdot9}}{6} = \dfrac{5 \mp \sqrt{52}}{3}
    \end{aligned}
\end{equation}
\begin{equation}
    \dfrac{2x}{x^2+1} = -\dfrac{1}{x} \qquad \qquad \text{C.E.: } x \ne 0 \quad \rightarrow \quad x \in \mathbb{R}\backslash \{0\}
\end{equation}
\begin{equation}
\begin{aligned}
    & 2x^2= - x^2 - 1 \\
    & 3x^2= - 1 \quad \rightarrow \quad \nexists x \in \mathbb{R}
\end{aligned}
\end{equation}

\subsection{Equazioni algebriche irrazionali}
\begin{definition}[Equazioni algebriche irrazionali] Le equazioni algebriche razionali sono equazioni che contengono polinomi, loro rapporti e potenze intere e non.
\end{definition}
\begin{example}[Esempi di equazioni algebriche irrazionali]
\begin{equation}
    \sqrt[3]{x-3} = 2 \qquad , \qquad \dfrac{2 x}{(x-1)^{\frac{1}{2}}} = - 2
\end{equation}
\end{example}
\subsubsection{Metodo di soluzione}
\begin{enumerate}
    \item Per prima cosa è necessario determinare le \textbf{condizioni di esistenza} di una soluzione. Bisogna richiedere che questa e tutte le operazioni scritte nel problema abbiano senso: bisogna richiedere che
        \begin{itemize}
            \item che non avvengano \textbf{divisioni} per zero;
            \item che siano non negativi i radicandi di eventuali \textbf{radici} con indice intero pari o non intero.
        \end{itemize}
    \item Successivamente, è possibile procedere con le semplificazioni per la ricerca della soluzione.
\end{enumerate}

\paragraph{Esempi}
\begin{equation}
    \sqrt[3]{x-3} = 2 \qquad \qquad \text{C.E.: } x \in \mathbb{R}
\end{equation}
\begin{equation}
    x-3 = 8 \qquad \rightarrow \qquad x = 11
\end{equation}
\begin{equation}
    \dfrac{2 x}{(x-1)^{\frac{1}{2}}} = - 2 \qquad \qquad \text{C.E.: } x-1>0 \quad \rightarrow \quad x \in (1,+\infty)
\end{equation}
\begin{equation}
\begin{aligned}
    &  x = -  (x-1)^{\frac{1}{2}} \\
    &  x^2 =   x-1  \\
    & x^2 - x + 1 = 0 \\
    & \Delta = (-1)^2 - 4 \cdot 1 \cdot 1 = -3 < 0 \qquad \rightarrow \qquad \nexists x \in \mathbb{R}
\end{aligned}
\end{equation}

\section{Equazioni non algebriche o trascendenti}
\subsection{Equazioni con i valori assoluti}
\subsection{Equazioni con esponenti e logaritmi}
\subsection{Equazioni con le funzioni armoniche}

\section{Metodi di soluzione approssimati}
\subsection{Metodo grafico}
\subsection{Metodi numerici}
{\color{red} Riferimento al capitolo dei metodi numerici}

% ==============================================================================
\chapter{Disequazioni}
\section{Disequazioni algebriche}
\section{Disequazioni non algebriche}

% ==============================================================================
\chapter{Sistemi di equazioni e di disequazioni}




\part{Algebra in $\mathbb{C}$}\label{book:complex_algebra}

% ==============================================================================
\chapter{Algebra complessa}\label{ch:complex-algebra}
\section{Definizione dei numeri complessi}
\begin{definition}[Unità immaginaria] 
 \begin{equation}
   i := \sqrt{-1}
 \end{equation}
\end{definition}

\begin{definition}[Numero complesso]
\begin{equation}
  z = x + i y \qquad , \qquad x,y \in \mathbb{R}
\end{equation}
\end{definition}


\subsection{Rappresentazione del piano complesso (di Argand--Gauss)}
Si può definire una relazione biunivoca tra l'insieme dei numeri complessi $\mathbb{C}$ e il piano $\mathbb{R}^2$.

\subsubsection{Rappresentazione cartesiana.}

\subsubsection{Rappresentazione polare.}
Trasformazione tra coordinate cartesiane e polari
\begin{equation}
\begin{cases}
  x = r \cos \theta \\
  y = r \sin \theta
\end{cases} \qquad , \qquad
\begin{cases}
  r = \sqrt{x^2 + y^2} \\
  \theta = \text{atan2}(x,y)
\end{cases}
\end{equation}
e quindi
\begin{equation}
  z = x + i y = r \left( \cos \theta + i \sin \theta \right)
\end{equation}

\paragraph{La relazione di Eulero e la rappresentazione polare dei numeri complessi.} Usando le espansioni in serie di Taylor delle funzioni $e^{i\theta}$, $\cos \theta$ e $\sin \theta$, Eulero ricavò la formula che da lui prende il nome
\begin{equation}
  e^{i \theta} = \cos \theta + i \sin \theta \ .
\end{equation}
\begin{proof}[Dimostrazione con le serie di Taylor]
\begin{equation}
\begin{array}{rll}
   \cos \theta  & = \sum_{n=0}^{\infty} (-1)^n \dfrac{\theta^{2n}}{(2n)!} & = 1 - \dfrac{\theta^2}{2} + \dfrac{\theta^4}{4!} + \dots \\
   \sin \theta  & = \sum_{n=0}^{\infty} (-1)^n \dfrac{\theta^{2n+1}}{(2n+1)!} & = \theta - \dfrac{\theta^3}{3!} + \dfrac{\theta^5}{5!} + \dots \\
   e^{i \theta} & = \sum_{n=0}^{\infty} \dfrac{(i\theta)^n}{n!} & = 1 + i\theta - \dfrac{\theta^2}{2} - i \dfrac{\theta^3}{3!} + \dfrac{\theta^4}{4!} + i \dfrac{\theta^5}{5!} +  \dots = \\
  & & = \left[  1 - \dfrac{\theta^2}{2} + \dfrac{\theta^4}{4!} + \dots \right] + i \left[ \theta - \dfrac{\theta^3}{3!} + \dfrac{\theta^5}{5!} + \dots \right]  = \\
  & & = \cos \theta + i \sin \theta \ .
\end{array}
\end{equation}
\end{proof}
%
\begin{proof}[Dimostrazione con la derivata] Si definisce la funzione con variabile reale $\theta \in \mathbb{R}$ e valori complessi $f: \mathbb{R} \rightarrow \mathbb{C}$, $f(\theta) = \cos \theta + i \sin \theta$, e si calcola la derivata rispetto alla variabile indipedente $\theta$
\begin{equation}
    \dfrac{d}{d \theta} \left( \cos \theta + i \sin \theta \right) = -\sin \theta + i \cos \theta = i \left( \cos \theta + i \sin \theta \right)
\end{equation}
    La funzione $f(\theta) = \cos \theta + i \sin \theta$ soddisfa quindi l'equazione differenziale
\begin{equation}
    \dfrac{d}{d \theta} f(\theta) = i f(\theta) \ ,
\end{equation}
    con condizione $f(0) = 1$. Essendo questo un problema lineare, se esiste, esiste una sola soluazione. Si può dimostrare che questa equazione è soddisfatta anche dalla funzione a variabile reale e valori complessi $g: \mathbb{R} \rightarrow \mathbb{C}$, $g(\theta) = e^{i\theta}$, la cui derivata è $g'(\theta) = i e^{i\theta} = i g(\theta)$.

    \noindent
    Possiamo quindi concludere l'uguaglianza tra le due funzioni $g(\theta) = f(\theta)$, $e^{i \theta} = \cos \theta + i \sin \theta$.
\end{proof}

% ------------------------------------------------------------------------------
\section{Operazioni con i numeri complessi}
In questa sezione vengono presentate le operazioni algebriche con i numeri complessi, osservando come convenga la forma cartesiana per somma e differenza, mentre la forma polare per prodotti, divisioni e soprattutto potenze e radici.

\subsection{Somma e differenza}
\begin{equation}
   z_1 + z_2 = x_1 + x_2 + i (y_1 + y_2) \ .
\end{equation}

\subsection{Prodotto e divisione}
\begin{equation}
   z_1 z_2 = r_1 r_2 e^{i(\theta_1+\theta_2)} \ .
\end{equation}

\subsection{Potenze e radici}
\begin{equation}
  z^n = \left( r e^{i\theta} \right)^n = r^n e^{i n\theta}
\end{equation}
\begin{equation}
  z^{\frac{1}{n}} = \left( r e^{i(\theta + 2\pi m)} \right)^{\frac{1}{n}} = r^{\frac{1}{n}} e^{i\left(\frac{\theta}{n} + \frac{m}{n} 2 \pi \right)}
\end{equation}

\subsection{Esponenziali e logaritmi}
\begin{equation}
\begin{aligned}
    e^z
    & = e^{x+i y} = e^x \, e^{i y} = e^x \left( \cos y + i \sin y \right) \\
    & = e^{r e^{i\theta}} = e^{r \left( \cos \theta + i \sin \theta \right)} = e^{r \cos \theta} e^{i \, r \, \sin \theta} = e^{r \cos \theta} \left[ \cos(r \sin \theta) + i \sin (r \sin \theta) \right]
\end{aligned}
\end{equation}
%
\begin{equation}
    \ln z = \ln \left( r e^{i \theta} \right) = \ln r + \ln e^{i \theta} = \ln r + i \, \theta
\end{equation}

% ==============================================================================
\chapter{Calcolo complesso}\label{ch:complex-calculus}


\part{Calcolo infinitesimale}


% =======================================================================================
\chapter{Funzioni $f:\mathbb{R} \rightarrow \mathbb{R}$ e limiti}

% ---------------------------------------------------------------------------------------
\section{Cenni di topologia per l'analisi}
\begin{definition}[Intervallo] Un intervallo è un sottoinsieme di $\mathbb{R}$, l'insieme dei numeri reali.
\end{definition}
\begin{definition}[Estremi superiore e inferiore] Se esiste, l'\textbf{estremo superiore} $M \in \mathbb{R}$ di un intervallo $E$ è il più piccolo maggiorante (cioè un numero uguale o maggiore di tutti i numeri dell'intervallo $E$), cioè:
    \begin{itemize}
        \item $M \ge x$ ,$\forall x \in E$
        \item $\nexists z$ maggiorante di $E$, con $z < M$
    \end{itemize}
    Se esiste, l'\textbf{estremo inferiore} $m \in \mathbb{R}$ di un intervallo $E$ è il più grande minorante (cioè un numero uguale o minore di tutti i numeri dell'intervallo $E$), cioè:
    \begin{itemize}
        \item $m \le x$ ,$\forall x \in E$
        \item $\nexists z$ minorante di $E$, con $z > m$
    \end{itemize}
\end{definition}
\begin{definition}[Intervallo aperto] Un intervallo aperto è un intervallo che non include nessuno dei suoi estremi.
\end{definition}
\begin{definition}[Intervallo chiuso] Un intervallo chiuso è un intervallo che include tutti i suoi estremi.
\end{definition}
Punti: accumulazione/isolato, interni/esterni/frontiera
Intervalli: limitati/illimitati, aperti/chiusi, intorno
\begin{definition}[Intorno di un punto] L'intorno completo $B(x_0, \varepsilon)$ del punto $x_0$ di raggio $\varepsilon$ è l'intervallo aperto
    \begin{equation}
        B(x_0, \varepsilon) = (x_0 - \varepsilon, x_0 + \varepsilon)  
    \end{equation}
\end{definition}
\begin{definition}[Intorno destro e sinisto] L'intorno destro di un punto $x_0$ della retta reale è l'insieme aperto
    \begin{equation}
        B^+(x_0,\varepsilon) = (x_0, x_0+\varepsilon) \ .
    \end{equation}
L'intorno destro di un punto $x_0$ della retta reale è l'insieme aperto
    \begin{equation}
        B^-(x_0,\varepsilon) = (x_0-\varepsilon, x_0) \ .
    \end{equation}
\end{definition}
\begin{definition}[Intorno di infinito]
L'intorno di $-\infty$ viene definito come l'insieme aperto
    \begin{equation}
        B(-\infty,M) = (-\infty,M)
    \end{equation}
L'intorno di $+\infty$ viene definito come l'insieme aperto
    \begin{equation}
        B(+\infty,M) = (M,+\infty)
    \end{equation}
\end{definition}

\begin{definition}[Punto di accumulazione] 
\end{definition}
\begin{definition}[Punto interno]
\end{definition}

\begin{definition}[{\color{red} Intervallo chiuso}] {\color{red} Un intervallo è chiuso se contiene tutti i suoi punti di accumulazione.}
\end{definition}
\begin{definition}[{\color{red} Intervallo aperto}] {\color{red} Un intervallo è aperto se ogni suo punto è un punto interno.}
\end{definition}

\begin{definition}[]
\end{definition}

% ---------------------------------------------------------------------------------------
\section{Funzioni}
\begin{definition}[Funzione a variabile e valori reali]
\end{definition}

% ---------------------------------------------------------------------------------------
\section{Limiti}
{\color{red}
\begin{itemize}
    \item Limite destro e limite sinistro
    \item Funzioni continue
\end{itemize}
}
\begin{definition}[Limite finito al finito] Il limite finito della funzione $f(x)$  per $x$ che tende a $x_0$,
    \begin{equation}
      \lim_{x \rightarrow x_0} f(x) = \ell
    \end{equation}
    è definito dalla condizione
    \begin{equation}
        \text{Per $\forall \varepsilon > 0$ tale che $|x-x_0| < \varepsilon$, $\exists \delta > 0$ tale che $|f(x) - \ell|< \delta$.}
    \end{equation}
\end{definition}
\begin{definition}[Limite infinito al finito] Il limite infinito della funzione $f(x)$  per $x$ che tende a $x_0$,
    \begin{equation}
      \lim_{x \rightarrow x_0} f(x) = \mp \infty
    \end{equation}
    è definito dalla condizione
    \begin{equation}
        \text{Per $\forall \varepsilon > 0$ tale che $|x-x_0| < \varepsilon$, $\exists M \lessgtr 0$ tale che $f(x) \lessgtr M$.}
    \end{equation}
\end{definition}
\begin{definition}[Limite finito al infinito] Il limite finito della funzione $f(x)$ per $x$ che tende all'infinito
    \begin{equation}
      \lim_{x \rightarrow \mp \infty} f(x) = \ell
    \end{equation}
    è definito dalla condizione
    \begin{equation}
        \text{Per $\forall N \lessgtr 0$ tale che $x \lessgtr  N$, $\exists \delta > 0$ tale che $|f(x) - \ell| < \delta$.}
    \end{equation}
\end{definition}
\begin{definition}[Limite infinito al infinito] Il limite infinito della funzione $f(x)$ per $x$ che tende all'infinito
    \begin{equation}
        \lim_{x \rightarrow \mp^{(1)} \infty} f(x) = \mp^{(2)} \infty
    \end{equation}
    è definito dalla condizione
    \begin{equation}
        \text{Per $\forall N \lessgtr^{(1)} 0$ tale che $x \lessgtr^{(1)} N$, $\exists M \lessgtr^{(2)} 0$ tale che $f(x) \lessgtr^{(2)} M$.}
    \end{equation}
\end{definition}

{\color{red}Grafici}

% ---------------------------------------------------------------------------------------
\section{Funzioni continue}
\begin{definition}[Funzione continua in un punto] Una funzione $f: \Omega \rightarrow \mathbb{R}$ è continua in un punto $x_0 \in \Omega$ se 
    \begin{equation}
        \lim_{x \rightarrow x_0}f(x) = f(x_0) \ .
    \end{equation}
\end{definition}
\begin{definition}[Funzione continua in un intervallo] Una funzione $f: \Omega \rightarrow \mathbb{R}$ è continua in un intervallo $[a, b] \subset \Omega$, se è una funzione continua in tutti i punti $x \in [a, b]$.
\end{definition}

\subsection{Teoremi sulle funzioni continue}
\begin{theorem}[Teorema di Weierstrass]\label{thm:weierstrass} Una funzione continua definita su un intervallo chiuso e limitato $[a, b]$ ammette in esso un massimo e un minimo assoluto.
\end{theorem}

% ---------------------------------------------------------------------------------------
\section{Teoremi sui limiti}
\begin{theorem}[Teorema del confronto]\label{thm:comparison}
\end{theorem}

\begin{theorem}[Teorema di de l'Hopital -- anticipazione]\label{thm:hopital:0} Il teorema di de l'Hopital riguarda il calcolo delle forme indeterminate $\frac{0}{0}$, $\frac{\infty}{\infty}$, cioè i limiti che possono essere scritti nella forma
    \begin{equation}
        \lim_{x \rightarrow x_0} \dfrac{f(x)}{g(x)}
    \end{equation}
    con
    \begin{itemize}
        \item $\lim_{x \rightarrow x_0} f(x) = \lim_{x \rightarrow x_0} g(x) = 0$ oppure
        \item $\lim_{x \rightarrow x_0} |f(x)| = \lim_{x \rightarrow x_0} |g(x)| = +\infty$
    \end{itemize}
\end{theorem}
    Per l'enunciato completo e la dimostrazione del teorema di de l'Hopital si rimanda al teorema \S\ref{thm:hopital} cap.\S\ref{ch:derivatives} sulle derivate.

% ---------------------------------------------------------------------------------------
\section{Infiniti e infinitesimi}

% ---------------------------------------------------------------------------------------
\section{Limiti notevoli}
\begin{center}
\begin{tabular}{ll}
    Potenza            & $\displaystyle\lim_{x \rightarrow 0} \dfrac{(1+x)^a - 1}{x} = a$ \\
    Seno               & $\displaystyle\lim_{x \rightarrow 0} \dfrac{\sin}{x} = 1$ \\
    Coseno             & $\displaystyle\lim_{x \rightarrow 0} \dfrac{1 - \cos x}{x^2} = \frac{1}{2}$ \\
    Esponenziale - def & $\displaystyle\lim_{x \rightarrow \infty} \left( 1 + \dfrac{1}{x} \right)^x = e$ \\
    \qquad " \qquad (1)& $\displaystyle\lim_{x \rightarrow 0} \dfrac{e^x - 1}{x} = 1$ \\
    \qquad " \qquad (2)& $\displaystyle\lim_{x \rightarrow 0} \dfrac{e^{x}}{1+x} = 1$ \\
    Logaritmo naturale & $\displaystyle\lim_{x \rightarrow 0} \dfrac{\ln (1+x)}{x} = 1$ \\
\end{tabular}
\end{center}

\subsection{Dimostrazioni}
\begin{proof}[Limite notevole del seno] Il limite viene calcolato usando il teorema del confronto (\ref{thm:comparison}), con le funzioni $x$, $\sin x$ e $\tan x$
    \begin{equation}
        \sin x \le x \le \tan x \qquad \rightarrow \qquad 1 \le \dfrac{x}{\sin{x}} \le \dfrac{1}{\cos x} \ .
    \end{equation}
    Per $x \rightarrow 0$, $\displaystyle\lim_{x \rightarrow 0} 1 = 1$ e $\displaystyle\lim_{x \rightarrow 0} \frac{1}{\cos x} = 1$ e quindi $ \lim_{x \rightarrow 0} \dfrac{ x }{\sin x} = 1 $.
\end{proof}
\begin{proof}[Limite notevole del coseno] Il limite viene calcolato usando il limite notevole del seno,
    \begin{equation}
        \begin{aligned}
            \lim_{x \rightarrow 0} \dfrac{1 - \cos x}{x^2} 
            & = \lim_{x \rightarrow 0} \dfrac{1 - \cos x}{x^2} \dfrac{1+\cos x}{1+\cos x} = \\
            & = \lim_{x \rightarrow 0} \dfrac{1 - \cos^2 x}{x^2} \dfrac{1}{1+\cos x} = \\ 
            & = \lim_{x \rightarrow 0} \underbrace{\dfrac{\sin^2 x}{x^2}}_{\left( \frac{\sin x}{x} \right)^2 \rightarrow 1} \underbrace{\dfrac{1}{1+\cos x}}_{\rightarrow \frac{1}{2}} = \frac{1}{2}
        \end{aligned}
    \end{equation}
\end{proof}

{\color{red}
\noindent
REMOVE from here, MOVE to chapter about series
\begin{definition}[Il numero di Nepero, $e$] Il numero di Nepero, $e$ può essere definito come il limite
    \begin{equation}
        e := \lim_{n \rightarrow \infty} \left( 1 + \dfrac{1}{n} \right)^n \ .
    \end{equation}
\end{definition}
Si prova l'esistenza al finito di questo limite, dimostrando che la serie $a_n := \left( 1 + \frac{1}{n} \right)^n $ è:
\begin{itemize}
    \item crescente in maniera monotona, i.e. $a_{n+1} > a_n$, $\forall n$
    \item limitata superiormente, i.e. $\exists M s.t. a_n < M$, $\forall n$
\end{itemize}
Per dimostrare che la serie è crescente, cerchiamo il legame tra due termini successivi.
Per dimostrare che la serie è limitata superiormente, dimostriamo che $a_n < 3$.
}

\begin{proof}[Limite notevole dell'esponenziale - definizione] Questo limite non è nient'altro che una delle definizioni del numero $e$, come mostrato nella Sez.\S\ref{sec:napier}.
\end{proof}
\begin{proof}[Limite notevole dell'esponenziale - (1)]
    \begin{equation}
        \begin{aligned}
            \lim_{x \rightarrow 0} \dfrac{e^x - 1}{x} 
            & = \lim_{x \rightarrow 0} \dfrac{1}{x}\sum_{k=0}^{\infty} \left( \dfrac{x^k}{k!} - 1 \right) = \\
            & = \lim_{x \rightarrow 0} \dfrac{1}{x}\sum_{k=1}^{\infty} \dfrac{x^k}{k!} = \\
            & = \lim_{x \rightarrow 0} \dfrac{1}{x} \left( x + \dfrac{x^2}{2} + \dfrac{x^3}{3!} + \dots  \right) = \\
            & = \lim_{x \rightarrow 0} \left( 1 + \dfrac{x}{2} + \dfrac{x^2}{3!} + \dots  \right) = 1 
        \end{aligned}
    \end{equation}
\end{proof}
\begin{proof}[Limite notevole dell'esponenziale - (2)]
    \begin{equation}
        \begin{aligned}
            \lim_{x \rightarrow 0} \dfrac{e^x}{1+x} 
            & = \lim_{x \rightarrow 0} \dfrac{1}{1+x}\sum_{k=0}^{\infty} \dfrac{x^k}{k!} = \\
            & = \lim_{x \rightarrow 0} \dfrac{1}{1+x}\left( 1 + x + \dfrac{x^2}{2} + \dfrac{x^3}{3!} + \dots \right) = \\
            & = \lim_{x \rightarrow 0} \bigg[ 1 + \underbrace{\dfrac{x^2}{1+x} \left( \dfrac{1}{2} + \dfrac{x}{3!} + \dots \right)}_{\rightarrow 0} \bigg] = 1 \\
        \end{aligned}
    \end{equation}
\end{proof}
\begin{proof}[Limite notevole del logaritmo]
    \begin{equation}
        \begin{aligned}
            \lim_{x \rightarrow 0} \dfrac{e^x}{1+x} 
            & = \lim_{x \rightarrow 0} \dfrac{1}{1+x}\sum_{k=0}^{\infty} \dfrac{x^k}{k!} = \\
            & = \lim_{x \rightarrow 0} \dfrac{1}{1+x}\left( 1 + x + \dfrac{x^2}{2} + \dfrac{x^3}{3!} + \dots \right) = \\
            & = \lim_{x \rightarrow 0} \bigg[ 1 + \underbrace{\dfrac{x^2}{1+x} \left( \dfrac{1}{2} + \dfrac{x}{3!} + \dots \right)}_{\rightarrow 0} \bigg] = 1 \\
        \end{aligned}
    \end{equation}
\end{proof}

% =======================================================================================
\chapter{Derivate}\label{ch:derivatives}

% ---------------------------------------------------------------------------------------
\section{Definizioni}
\begin{definition}[Rapporto incrementale]
\begin{equation}
    \dfrac{f(x+\Delta x) - f(x)}{\Delta x}
\end{equation}
\end{definition}

\begin{definition}[Derivata]
\begin{equation}
  f'(x) = \dfrac{d }{dx} f(x) := \lim_{\Delta x \rightarrow 0} \dfrac{f(x+\Delta x) - f(x)}{\Delta x}
\end{equation}
\end{definition}

\begin{definition}
\end{definition}

\paragraph{Interpretazione geometrica.}

% ---------------------------------------------------------------------------------------
\section{Regole di derivazione}
\subsection{Regole}
\paragraph{Derivata della somma di due funzioni e il prodotto per uno scalare}
\begin{equation}
\begin{aligned}
    (f(x) + g(x))' & = f'(x) + g'(x) \\
    (a f(x))' & = a f'(x) \\
\end{aligned}
\end{equation}
\begin{property}[Operatore lineare] La derivata è un operatore lineare.
\end{property}

\paragraph{Derivata del prodotto di due funzioni}
\begin{equation}
    (f(x) g(x))' = f'(x) g(x) + f(x) g'(x)
\end{equation}
\paragraph{Derivata del rapporto di due funzioni}
\begin{equation}
    \left( \dfrac{f(x)}{g(x)} \right)' = \dfrac{f'(x)g(x) - f(x)g'(x)}{g^2(x)}
\end{equation}
\paragraph{Derivata di una funzione composta}
\begin{equation}
    \dfrac{d}{dx} f(g(x)) = \dfrac{d}{dy} f(y) \bigg|_{y=g(x)} \dfrac{d}{dx} g(x)
\end{equation}
\paragraph{Derivata della funzione inversa}
\begin{equation}
\end{equation}

\subsection{Dimostrazioni}

\paragraph{Derivata della somma di due funzioni}
\paragraph{Derivata del prodotto di due funzioni}
\begin{equation}
\begin{aligned}
\dfrac{d}{dx} \left(f(x) g(x) \right) & =
     \lim_{\Delta x \rightarrow 0} \dfrac{f(x+\Delta x) g(x+\Delta x) - f(x)g(x) }{\Delta x} =  \\
 & = \lim_{\Delta x \rightarrow 0} \dfrac{f(x+\Delta x) g(x+\Delta x) - f(x+\Delta x) g(x) + f(x+\Delta x) g(x) - f(x)g(x) }{\Delta x} =  \\
 & = \lim_{\Delta x \rightarrow 0} \dfrac{f(x+\Delta x) g(x+\Delta x) - f(x+\Delta x) g(x) + f(x+\Delta x) g(x) - f(x)g(x) }{\Delta x} =  \\
    & = \lim_{\Delta x \rightarrow 0} f(x+\Delta x) \dfrac{ g(x+\Delta x) - g(x)}{\Delta x} + \lim_{\Delta x \rightarrow 0} \dfrac{ f(x+\Delta x)  - f(x) }{\Delta x} g(x) =  f'(x)g(x) + f(x)g'(x) \ .
\end{aligned}
\end{equation}
\paragraph{Derivata del rapporto di due funzioni}
\begin{equation}
\begin{aligned}
    \dfrac{d}{dx} \left(\dfrac{f(x)}{ g(x)} \right) & =
     \lim_{\Delta x \rightarrow 0} \dfrac{1}{\Delta x} \left[ \dfrac{f(x+\Delta x)}{ g(x+\Delta x)} - \dfrac{f(x)}{g(x)} \right] =  \\
    & = \lim_{\Delta x \rightarrow 0} \dfrac{1}{\Delta x} \dfrac{f(x+\Delta x) g(x) - f(x)g(x+\Delta x)}{g(x+\Delta x)g(x)}  =  \\
    & = \lim_{\Delta x \rightarrow 0} \dfrac{1}{g(x+\Delta x)g(x)} \dfrac{f(x+\Delta x) g(x) - f(x) g(x) + f(x)g(x)  - f(x)g(x+\Delta x)}{\Delta x}  =  \\
    & = \lim_{\Delta x \rightarrow 0} \dfrac{1}{g(x+\Delta x)g(x)} \dfrac{f(x+\Delta x) - f(x)}{\Delta x} g(x) - \lim_{\Delta x \rightarrow 0} \dfrac{1}{g(x+\Delta x)g(x)} \dfrac{g(x+\Delta x) - g(x)}{\Delta x} f(x) =  \\
    & = \dfrac{f'(x) g(x)}{g^2(x)} - \dfrac{f(x)g'(x)}{g^2(x)} = \dfrac{f'(x)g(x) - f(x) g'(x)}{g^2(x)}
\end{aligned}
\end{equation}
\paragraph{Derivata di una funzione composta}
\begin{equation}
\begin{aligned}
    \dfrac{d}{dx} f(g(x)) & =
     \lim_{\Delta x \rightarrow 0} \dfrac{1}{\Delta x} \left[ f(g(x+\Delta x)) - f(g(x)) \right] =  \\
    & = \lim_{\Delta x \rightarrow 0} \dfrac{f(g(x+\Delta x)) - f(g(x))}{g(x+\Delta x) - g(x)} \dfrac{g(x+\Delta x) - g(x)}{\Delta x} =  \\
    & = \dots \\
    & = f'(g(x)) \, g'(x) \ .
\end{aligned}
\end{equation}

% ---------------------------------------------------------------------------------------
\section{Teoremi}
\begin{theorem}[Teorema di Fermat]\label{thm:fermat} Sia $f: \ (a,b) \rightarrow \mathbb{R}$ una funzione definita sull'intervallo $(a,b)$ e derivabile nel punto $x_0 \in (a,b)$, punto di estremo locale. Allora $f'(x_0) = 0$.
\end{theorem}

\begin{theorem}[Teorema di Rolle]\label{thm:rolle} Sia $f: \ [a,b] \rightarrow \mathbb{R}$ una funzione continua e derivabile in ogni punto dell'intervallo $[a,b]$ che assume valori uguali $f(a) = f(b)$, allora esiste almeno un punto $c \in (a,b)$ in cui la derivata della funzione si annulla, $f'(c) = 0$.
\end{theorem}
\begin{proof}
    Secondo il teorema di Weierstrass (\ref{thm:weierstrass}) sulle funzioni continue, la funzione $f$ ammette punti $c \in [a,b]$ di massimo e un minimo globali. Si possono riconoscere ora due casi:
    \begin{itemize}
        \item se i punti di massimo e minimo globali coincidono entrambi con gli estremi dell'intervallo, si può scrivere
            \begin{equation} m = f(a) = f(b) \le f(x) \le M = f(a) = f(b) \end{equation}
                e quindi la funzione è costante, $f(x) = f(a) = f(b)$, e la sua derivata è nulla in tutti i punti dell'intervallo,$f'(x) = 0, \ \forall x \in (a,b)$;
            \item altrimenti, esiste un punto di massimo o minimo globale $x=c \in (a,b)$ interno all'intervallo; i punti di massimo o minimo globali globali sono anche punti di massimo o minimo locali e quindi, per il teorema di Fermat (\ref{thm:fermat}), $f'(c) = 0$.
    \end{itemize} 
\end{proof}

\begin{theorem}[Teorema di Cauchy]\label{thm:cauchy} Siano $f, g: \ [a,b] \rightarrow \mathbb{R}$ due funzioni continue in $[a,b]$ e derivabili su $(a,b)$, allora esiste un punto $c \in (a,b)$ tale che
    \begin{equation}
        \left[ f(b) - f(a)  \right]g'(c) = \left[ g(b) - g(a)  \right]f'(c) \ .
    \end{equation}
\end{theorem}
\begin{proof}
  La dimostrazione del teorema di Cauchy può essere svolta con un'opportuna scelta della funzione alla quale applicare il teorema di Rolle (\ref{thm:rolle}). Ad esempio, la funzione
    \begin{equation}
       h(x) = [f(b) - f(a)] [ g(x) - g(a)] - [g(b)-g(a)] [f(x) - f(a)] \ ,
    \end{equation}
    soddisfa le ipotesi del teoremda di Rolle, poichè
    \begin{itemize}
        \item $h(x)$ è continua su $[a,b]$ e derivabile su $(a,b)$
        \item $h(a) = h(b)$, dal calcolo diretto
            \begin{equation}
                \begin{aligned}
                    h(a) & = [f(b) - f(a)] [ g(a) - g(a)] - [g(b)-g(a)] [f(a) - f(a)] = 0 \\ 
                    h(b) & = [f(b) - f(a)] [ g(b) - g(a)] - [g(b)-g(a)] [f(b) - f(a)] = 0  \ .
                \end{aligned}
            \end{equation}
    \end{itemize}
    Applicando il teorema di Rolle alla funzione $h(x)$, possiamo concludere che esiste un punto $c \in (a,b)$ tale che $h'(c) = 0$ e quindi
    \begin{equation}
        0 = h'(c) = [f(b) - f(a)] g'(x) - [g(b)-g(a)] f'(x) \ .
    \end{equation}
\end{proof}

\begin{theorem}[Teorema di Lagrange]\label{thm:lagrange} Sia $f: \ [a,b] \rightarrow \mathbb{R}$ una funzione continua in $[a,b]$ e derivabile in $(a,b)$, allora esiste un punto $c \in (a,b)$ tale che
    \begin{equation}
        f(b) - f(a) = f'(c) (b-a) \ .
    \end{equation}
\end{theorem}
\begin{proof}
    La dimostrazione del teorema di Lagrange segue direttamente dalla dimostrazione del teorema di Cauchy (\ref{thm:cauchy}), usando le funzioni $f(x)$ e $g(x) = x$, la cui derivata è $g'(x) = 1$.
\end{proof}

\subsection{Teorema di de l'Hopital}
\begin{theorem}[Teorema di de l'Hopital]\label{thm:hopital}
\end{theorem}

% ------------------------------------------------------------------------------
\section{Derivate fondamentali}
\begin{center}
    \begin{tabular}{lll}
        Potenza      & $f(x) = x^a$       & $f'(x) = a x^{a-1}$    \\
        Esponenziale & $f(x) = e^x$       & $f'(x) = e^x$          \\
        Logaritmo    & $f(x) = \ln x$     & $f'(x) = \frac{1}{x}$  \\
        Seno         & $f(x) = \sin x$    & $f'(x) = \cos x$       \\
        Coseno       & $f(x) = \cos x$    & $f'(x) =-\sin x$       \\
    \end{tabular}
\end{center}
\subsection{Dimostrazioni}
\begin{proof}[Derivata della potenza]
\begin{equation}
    \begin{aligned}
    (x^n)'
        & = \lim_{\varepsilon \rightarrow 0 } \dfrac{(x+\varepsilon)^n - x^n}{\varepsilon} = \\
        & = \lim_{\varepsilon \rightarrow 0 } \dfrac{ \left( 1 + \frac{\varepsilon}{x} \right)^n - 1 }{\varepsilon} x^n = \\
        & = \underbrace{\lim_{\varepsilon \rightarrow 0 } \dfrac{ \left( 1 + \frac{\varepsilon}{x} \right)^n - 1 }{\frac{\varepsilon}{x}}}_{=n} x^{n-1} = \\
        & = n x^{n-1}
    \end{aligned}
\end{equation}
\end{proof}

\begin{proof}[Derivata della funzione esponenziale]
\begin{equation}
    \begin{aligned}
        (e^x)' 
        & = \lim_{\varepsilon \rightarrow 0 } \dfrac{e^{x+\varepsilon} - e^x}{\varepsilon} = \\
        & = \underbrace{\lim_{\varepsilon \rightarrow 0 } \dfrac{(e^{\varepsilon} - 1)}{\varepsilon}}_{= 1} e^{x} = \\
        & = e^x
    \end{aligned}
\end{equation}
\end{proof}
\begin{proof}[Derivata della funzione logaritmo] {\color{red} Il dominio della funzione logaritmo è $\mathbb{R}^+$, ossia l'insieme dei numeri reali positivi. Essendo l'argomento del logaritmo diverso da zero, si può dividere per esso senza il timore di incorrere in espressioni non definite}
\begin{equation}
    \begin{aligned}
        (\ln x)' 
        & = \lim_{\varepsilon \rightarrow 0 } \dfrac{\ln(x+\varepsilon) - \ln x}{\varepsilon} = \\
        & = \lim_{\varepsilon \rightarrow 0 } \dfrac{\ln \left( x\left( 1 + \frac{\varepsilon}{x} \right) \right) - \ln x}{\varepsilon} = \\
        & = \lim_{\varepsilon \rightarrow 0 } \dfrac{\ln  x + \ln \left( 1 + \frac{\varepsilon}{x}  \right) - \ln x}{\varepsilon} = \\
        & = \lim_{\varepsilon \rightarrow 0 } \dfrac{ \ln \left( 1 + \frac{\varepsilon}{x}  \right) }{\varepsilon} = \\
        & = \underbrace{\lim_{ \small\begin{matrix} \varepsilon \rightarrow 0 \\  \varepsilon / x \rightarrow 0  \end{matrix}} \dfrac{ \ln \left( 1 + \frac{\varepsilon}{x}  \right) }{\frac{\varepsilon}{x}}}_{=1} \dfrac{1}{x} = \\
        & = \dfrac{1}{x}
    \end{aligned}
\end{equation}
\end{proof}
\begin{proof}[Derivata della funzione seno]
\begin{equation}
    \begin{aligned}
        (\sin x)' 
        & = \lim_{\varepsilon \rightarrow 0 } \dfrac{\sin(x+\varepsilon) - \sin x}{\varepsilon} = \\
        & = \lim_{\varepsilon \rightarrow 0 } \dfrac{\sin x \cos\varepsilon + \sin \varepsilon \cos x - \sin x}{\varepsilon} = \\
        & = \underbrace{\lim_{\varepsilon \rightarrow 0 } \left( \dfrac{ \cos\varepsilon - 1}{\varepsilon} \right) }_{\lim_{\varepsilon \rightarrow 0} \left(-\frac{1}{2} \varepsilon \right) = 0 } \sin x +  \underbrace{\lim_{\varepsilon \rightarrow 0 } \dfrac{\sin \varepsilon}{\varepsilon} }_{= 1} \cos x  = \\
        & = \cos x 
    \end{aligned}
\end{equation}
\end{proof}
\begin{proof}[Derivata della funzione coseno]
\begin{equation}
    \begin{aligned}
        (\cos x)' 
        & = \lim_{\varepsilon \rightarrow 0 } \dfrac{\cos(x+\varepsilon) - \cos x}{\varepsilon} = \\
        & = \lim_{\varepsilon \rightarrow 0 } \dfrac{\cos x \cos\varepsilon - \sin \varepsilon \sin x - \cos x}{\varepsilon} = \\
        & = \underbrace{\lim_{\varepsilon \rightarrow 0 } \left( \dfrac{ \cos\varepsilon - 1}{\varepsilon} \right) }_{\lim_{\varepsilon \rightarrow 0} \left(-\frac{1}{2} \varepsilon \right) = 0 } \cos x -  \underbrace{\lim_{\varepsilon \rightarrow 0 } \dfrac{\sin \varepsilon}{\varepsilon} }_{= 1} \sin x  = \\
        & = - \sin x 
    \end{aligned}
\end{equation}
\end{proof}

\section{Derivate di ordine superiore}

\section{Espansioni in serie di Taylor e MacLaurin}\label{ch:taylor}
\begin{definition}[Serie di Taylor] La serie di Taylor di una funzione $f(x)$ centrata in $x=x_0$ è la serie polinomiale
\begin{equation}
   T[f(x_0)](x) := \sum_{n=0}^{\infty} \dfrac{f^{(n)}(x_0)}{n!} (x-x_0)^n \ .
\end{equation}
\end{definition}
%
\begin{theorem}
La serie di Taylor troncata alla $n$-esima potenza,
\begin{equation}
  T^n[f(x_0)](x) = \sum_{k=0}^{n} \dfrac{f^{(k)}(x_0)}{i!} (x-x_0)^k \ ,
\end{equation}
è un'approssimazione dell'$n$-esimo ordine della funzione $f(x)$, i.e.
\begin{equation}
  f(x) - T^n[f(x_0)](x) \sim o(|x-x_0|^{n})
\end{equation}
\end{theorem}

\begin{definition}[Serie di MacLaurin] La serie di MacLaurin di una funzione $f(x)$ è definita come la sua serie di Taylor centrata in $x=0$.
\end{definition}


\section{Applicazioni}
\subsection{Punti in cui $f'(x) = 0$: punti di estremo locale e flessi orizzontali}
Per una funzione derivabile, i punti in cui si annulla la derivata prima di una funzione $f(x)$ possono essere:
\begin{itemize}
    \item punti di massimo locale
    \item punti di minimo locale
    \item punti di flesso orizzontale
\end{itemize}
\begin{definition}[Punto di massimo locale]
\end{definition}
\begin{theorem}[Condizioni per un punto di massimo locale] Il punto $x_0$ è un punto di massimo locale per una funzione $f(x)$ derivabile se
    \begin{itemize}
        \item $f'(x_0) = 0$
        \item la prima derivata di ordine superiore valutata in $x_0$ diversa da zero, è una derivata di \textbf{ordine pari} ed è \textbf{negativa}, $f^{(n)}(x_0) < 0$
    \end{itemize}
\end{theorem}
\begin{definition}[Punto di minimo locale]
\end{definition}
\begin{theorem}[Condizioni per un punto di minimo locale] Il punto $x_0$ è un punto di mimimo locale per una funzione $f(x)$ derivabile se
    \begin{itemize}
        \item $f'(x_0) = 0$
        \item la prima derivata di ordine superiore valutata in $x_0$ diversa da zero, è una derivata di \textbf{ordine pari}, ed è \textbf{positiva}, $f^{(n)}(x_0) > 0$
    \end{itemize}
\end{theorem}
\begin{definition}[Punto di flesso orizzontale]
\end{definition}
\begin{theorem}[Condizioni per un punto di flesso orizzontale] Il punto $x_0$ è un punto di flesso orizzontale per una funzione $f(x)$ derivabile se
    \begin{itemize}
        \item $f'(x_0) = 0$
        \item la prima derivata di ordine superiore valutata in $x_0$ diversa da zero, è una derivata di \textbf{ordine dispari}
    \end{itemize}
\end{theorem}

\subsection{Punti in cui $f''(x) = 0$: punti di flesso}

\subsection{Studio funzione}

% =======================================================================================
\chapter{Integrali}\label{ch:integrals}

\section{Definizioni}
\begin{definition}[Somma di Riemann] Data una funzione continua $f:[a,b] \rightarrow \mathbb{R}$, e una partizione $\mathcal{P} = \left\{x_0, x_1, \dots, x_n | a = x_0 < x_1 < \dots x_n = b \right\}$ dell'intervallo $[a,b]$, la somma di Riemann viene definita come
    \begin{equation}
        \sigma_n = \sum_{i=1}^{n} f(\xi_i) \ (x_{i} - x_{i-1})
    \end{equation}
con $\xi_i \in [x_{i-1}, x_i]$.
\end{definition}

\begin{definition}[Integrale di Riemann]\label{def:riemann-integral} Definendo $\Delta x := \max_i(x_i - x_{i-1}) $, l'integrale di Riemann viene definito come il limite della somma di Riemann per $\Delta x  \rightarrow 0$ (e di conseguenza il numero di intervalli della partizione $n \rightarrow \infty$), e viene indicato come
    \begin{equation}
        \int_{x=a}^b f(x) dx = \lim_{\Delta x \rightarrow 0} \sigma_n
    \end{equation}
\end{definition}

\begin{definition}[Integrale definito] L'integrale della definizione (\ref{def:riemann-integral}) svolto tra due valori di $x$ viene definito integrale definito. 
\end{definition}

\paragraph{Interpretazione geometrica}

\begin{definition}[Integrale indefinito]
\end{definition}

\section{Proprietà}
\begin{equation}
    \int_{x=a}^{b} \left[ A f(x) + B g(x) \right] dx = A \int_{x=a}^{b} f(x) dx + B \int_{x=a}^{b} g(x) dx
\end{equation}
\begin{equation}
    \int_{x=a}^{b} f(x) dx = \int_{x=a}^{c} f(x) dx + \int_{x=c}^{b} f(x) dx
\end{equation}

\section{Teoremi}
\begin{definition}[Media]
    \begin{equation}
        M(f, [a,b]) = \dfrac{1}{b-a} \int_{x=a}^{b} f(x) dx
    \end{equation}
\end{definition}

\begin{theorem}[Teorema della media]
    Sia $f: \ [a,b] \rightarrow \mathbb{R}$ una funzione continua su $[a,b]$, allora esiste un punto $c \in [a,b]$ tale che
    \begin{equation}
        \int_{x=a}^{b} f(x) dx = f(c) (b-a) \ .
    \end{equation}
\end{theorem}
\begin{proof}
    La dimostrazione è immediata, applicando il teorema di Lagrange (\ref{thm:lagrange}) alla primitiva $F(x)$ della funzione $f(x)$, $F(x) = \int_{t=x_0}^{x} f(t) dt$ 
\begin{equation}
    \int_{t=a}^{b} f(x) dx = F(b) - F(a) = F'(c) (b-a) = f(c) (b-a) \ .
\end{equation}
\end{proof}


\begin{theorem}[Teorema fondamentale del calcolo infinitesimale]
    \begin{equation}
        \dfrac{d}{dx} \int_{t=a}^{x} f(t) dt = f(x)
    \end{equation}
\end{theorem}
\begin{equation}
    \begin{aligned}
        \dfrac{d}{dx} \int_{t=a}^{x} f(t) dt & = 
        \lim_{\Delta x \rightarrow 0} \dfrac{1}{\Delta x} \left[ \int_{t=a}^{x+\Delta x} f(t) dt - \int_{t=a}^{x} f(t) dt\right] = \\
        & = \lim_{\Delta x \rightarrow 0} \dfrac{1}{\Delta x} \int_{t=x}^{x+\Delta x} f(t) dt = \\
        & = \lim_{\Delta x \rightarrow 0} \dfrac{1}{\Delta x} \Delta x \, f(\xi) = \qquad \qquad \qquad \text{(con $\xi \in [x, x+\Delta x]$)} \\
        & = f(x)  \  .
    \end{aligned}
\end{equation}

\section{Integrali fondamentali}

\section{Regole di integrazione}
\subsection{Integrazione per parti}
\begin{itemize}
 \item Definendo $F(x)$, $G(x)$ le primitive delle funzioni $f(x)$, e $g(x)$
 \item Integrando in $x$ dalla regola di \textbf{derivazione del prodotto} $(F(x)G(x))' = F'(x)G(x) + F(x)G'(x)$, riscritta isolando il termine $F'(x)G(x) = (F(x)G(x))' - F(x)G'(x)$
\end{itemize}
si ottiene
\begin{equation}
\begin{aligned}
    \int f(x) G(x) dx & = \int (F(x) G(x))' dx - \int F(x) G'(x) dx = \\
    &= F(x)G(x) - \int F(x) G'(x) dx 
\end{aligned}
\end{equation}

\begin{example}
\end{example}

\subsection{Integrazione con sostituzione}
\begin{itemize}
    \item Definendo la funzione composta $\overline{F}(x) = F(y(x))$ e le derivate
        \begin{equation}
            \overline{f}(x) = \dfrac{d}{dx} \overline{F}(x) \qquad , \qquad f(y) = \dfrac{d}{dy}F(y)
        \end{equation}
    \item Partendo dalla regola di \textbf{derivazione della funzione composta}, $\overline{F}(x) = F(y(x))$
        \begin{equation}
            \overline{f}(x) = \dfrac{d}{dx} \overline{F}(x) = \dfrac{d}{dx} F(y(x)) = \dfrac{d F}{dy}(y(x)) \dfrac{d y}{d x}(x) = f(y(x)) y'(x)
        \end{equation}
\end{itemize}
Usando il teorema fondamentale del calcolo infinitesimale
\begin{equation}
\begin{aligned}
    F(y) & = \int f(y) dy \\
    \overline{F}(x) & = \int \overline{f}(x) dx = \int f(y(x)) \ y'(x) dx
\end{aligned}
\end{equation}
Se si introduce la dipendenza $y(x)$ nella prima equazione, si ottiene l'uguaglianza tra le ultime due espressioni, $F(y(x)) = \overline{F}(x)$, e quindi
\begin{equation}
  \int f(y) dy = \int f(y(x)) \ y'(x) dx \ .
\end{equation}

\begin{example}
\end{example}


\part{Equazioni differenziali ordinarie}

% ==============================================================================
\chapter{Introduzione}
\begin{definition}[Equazione differenziale ordinaria] Un'equazione differenziale ordinaria è un'equazione che ha come incognita una funzione $y(x)$, nella quale possono comparire la funzione incognita $y(x)$, le sue derivate $y^{(n)}(x)$ e la variabile indipendente $x$, che può essere scritto nella forma implicita
\begin{equation}\label{ode:def}
  F\left(x,y(x),y'(x), y''(x), \dots y^{(n)}(x) \right) = 0 \ , \qquad \text{con $x \in \Omega = [a,b]$}.
\end{equation}
L'\textbf{ordine} dell'equazione differenziale viene definito come l'ordine massimo delle derivate della funzione incognita che compaiono nell'equazione.
\end{definition}

\begin{definition}[Equazione differenziale ordinaria lineare] Un'equazione differenziale è lineare se si può scrivere come l'uguaglianza di una combinazione lineare delle derivate della funzione incognita e una funzione nota, $f(x)$. Ad esempio, la forma generale dell'equazione differenziale ordinaria di ordine $n$ può essere scritta come
\begin{equation}
    a_n(x) y^{(n)}(x) + a_{n-1}(x) y^{(n-1)}(x) + \dots + a_1(x) y'(x) + a_0(x) y(x) = f(x) \ , \qquad \text{con $x \in \Omega$}.
\end{equation}
\end{definition}

\begin{definition}[Equazione differenziale ordinaria lineare a coefficienti costanti] Un'equazione differenziale lineare a coefficienti costanti è un'equazione differenziale ordinaria lineare con coefficienti $a_i(x) = a_i$, numeri che non dipendono dalla variabile indipendente $x$,
\begin{equation}
    a_n y^{(n)}(x) + a_{n-1} y^{(n-1)}(x) + \dots + a_1 y'(x) + a_0 y(x) = f(x) \ , \qquad \text{con $x \in \Omega$}.
\end{equation}
\end{definition}
%
\begin{definition}[Equazione differenziale ordinaria lineare omogenea a coefficienti costanti] \color{red}{Un'equazione differenziale lineare omogenea a coefficienti costanti è un'equazione differenziale ordinaria lineare a coefficienti costanti con $f(x) = 0$},
\begin{equation}
    a_n y^{(n)}(x) + a_{n-1} y^{(n-1)}(x) + \dots + a_1 y'(x) + a_0 y(x) = 0 \ , \qquad \text{con $x \in \Omega$}.
\end{equation}
\end{definition}
%
In generale, la soluzione dell'equazione (\ref{ode:def}) dipende da $n$ parametri indeterminati. In generale, un problema differenziale è composto da:
\begin{itemize}
    \item un'equazione differenziale di ordine $n$
    \item $n$ condizioni per determinare gli $n$ parametri altrimenti indeterminati
\end{itemize}

\begin{definition}[Problema di Cauchy] Un problema di Cauchy è definito da:
\begin{itemize}
  \item un'equazione differenziale di ordine $n$
  \begin{equation}
     F\left(x,y(x),y'(x), y''(x), \dots y^{(n)}(x) \right) = 0 \ , \qquad \text{con $x \in \Omega = [a,b]$}.
  \end{equation}
  \item $n$ condizioni che definiscono il valore della funzione incognita e delle prime $n-1$ derivate nell'estremo inferiore dell'intervallo
  \begin{equation}
  \begin{aligned}
    y(a) & = y_0 \\
    y'(a) & = y_1 \\
    \dots \\
    y^{(n-1)}(a) & = y_{n-1} \\
  \end{aligned}
  \end{equation}
\end{itemize}

\end{definition}

% ==============================================================================
\chapter{Equazioni differenziali ordinarie lineari a coefficienti costanti}

\section{Equazioni differenziali lineari a coefficienti costanti di primo ordine}
\subsection{Equazioni differenziali lineari omogenee a coefficienti costanti di primo ordine}
\begin{equation}\label{ode:first_order}
  a y'(x) + b y(x) = 0 \ , \qquad \text{con $x \in \Omega$ e  $a \ne 0$}
\end{equation}
Si cerca la soluzione nella forma $y(x) = \alpha e^{\beta x}$ e, calcolando la derivata e inserendo nell'equazione, si ottiene
\begin{equation}
  ( a \beta + b ) \alpha e^{\beta x} = 0 \ .
\end{equation}
Il prodotto di tre fattori si annulla quando si annulla uno dei tre fattori:
\begin{itemize}
    \item $e^{\beta x}$ non si annulla per nessun valore di $x$
    \item se si annulla $\alpha$, $\alpha = 0$, si otterebbe la soluzione triviale $y(x) = 0$
    \item $\rightarrow$ deve quindi annullarsi il fattore $a \beta + b$: si ottiene quindi il valore $\beta = -\dfrac{b}{a}$
\end{itemize}
La forma generale della soluzione dell'equazione (\ref{ode:first_order}) è quindi
\begin{equation}
  y(x) = \alpha e^{-\frac{b}{a}x}
\end{equation}
Per determinare il coefficiente $\alpha$ è necessaria una condizione che definisca il valore della funzione (o della sua derivata) in un punto del dominio o del suo contorno. 

\section{Equazioni differenziali lineari a coefficienti costanti di secondo ordine}
\subsection{Equazioni differenziali lineari omogenee a coefficienti costanti di primo ordine}
\begin{equation}
  a y''(x) + b y'(x) + c y(x) = 0 \ , \qquad \text{con $x \in \Omega$ e  $a \ne 0$}
\end{equation}
Si cerca la soluzione nella forma $y(x) = \alpha e^{\beta x}$ e, calcolando le derivate e inserendo nell'equazione, si ottiene
\begin{equation}
  ( a \beta^2 + b \beta + c ) \alpha e^{\beta x} = 0 \ .
\end{equation}
I valori di $\beta$ si ottiengono dalla soluzione dell'equazione di secondo grado in $\beta$, $a \beta^2 + b \beta + c = 0$ che, a seconda del segno del discriminante $\Delta = b^2 - 4ac$, possono essere:
\begin{itemize}
    \item $\Delta > 0$: esistono due soluzioni reali distinte $\beta_{1,2} = \dfrac{-b\mp\sqrt{\Delta}}{2a}$.
    \item $\Delta = 0$: esistono due soluzioni reali coincidenti $\beta_{1,2} = -\dfrac{b}{2a}$
    \item $\Delta < 0$: esistono due soluzioni complesse coniugate $\beta_{1,2} = \dfrac{-b}{2a} \mp j \dfrac{\sqrt{-\Delta}}{2a}$
\end{itemize}
La soluzione dell'equazione differenziale assume quindi la forma
\begin{itemize}
    \item $\Delta > 0$:
    \begin{equation}
       y(x) = \alpha_1 e^{\beta_1 x} + \alpha_2 e^{\beta_2 x}
    \end{equation}
    \item $\Delta = 0$:
    \begin{equation}
       y(x) = \alpha_1 e^{\beta x} + \alpha_2 x e^{\beta x}
    \end{equation}
    \item $\Delta < 0$:
    \begin{equation}
    \begin{aligned}
       y(x) & = \alpha_1 e^{\beta x} + \alpha_2 e^{\beta^* x} = \\
            & = \alpha_1 e^{(re\{ \beta \} + i im\{\beta\})x} + \alpha_2 e^{(re\{ \beta \} - i im\{\beta\}) x} = \\
            & = e^{re\{ \beta \}x} \left( \alpha_1 e^{i \ im\{\beta\} x} +  \alpha_2 e^{-i \ im\{\beta\}x }  \right) \ ,
    \end{aligned}
    \end{equation}
    e per avere una soluzione reale, bisogna imporre $\alpha_2 = \alpha_1^*$, per ottenere la somma di due numeri complessi coniugati, uguale al doppio della somma della loro parte reale,
    \begin{equation}
    \begin{aligned}
       y(x)  & = 2 e^{re\{ \beta \}x} \left( re\{ \alpha_1 \} \cos(\beta x) - im\{\alpha_1 \} \sin(\beta x) \right) \ ,
    \end{aligned}
    \end{equation}
che può essere riscritta come
    \begin{equation}
    \begin{aligned}
       y(x)  & = e^{re\{ \beta \}x} \left( A \cos(\beta x) + B \sin(\beta x) \right) = \\ 
             & = C e^{re\{ \beta \}x} \cos(\beta x + \phi)  \ .
    \end{aligned}
    \end{equation}
\end{itemize}

% ==============================================================================
\chapter{Metodo di separazione delle variabili}



\part{Vettori}
% Introduzione: motivazioni
Motivazione:
\begin{itemize}
  \item non tutti gli oggetti di interesse della Matematica, della Fisica o delle Scienze in generale possono essere adeguatamente rappresentati da un singolo numero
  \item esempi: posizione, velocità, forza,\dots
\end{itemize}

Storia:
\begin{itemize}
  \item \dots
  \item da vettori nello spazio fisico a struttura astratta matematica
\end{itemize}


\chapter{Algebra vettoriale}

\section{Definizioni}
\begin{definition}[Spazio vettoriale] Uno spazio vettoriale è una struttura matematica composta da:
    \begin{itemize}
        \item un insieme $V$, i cui elementi $\mathbf{v} \in V$ sono chiamati \textbf{vettori}
        \item un campo $F$, i cui elementi $a \in F$ sono chiamati \textbf{scalari}
        \item due operazioni chiuse rispetto a $V$, cioè il cui risultato è un elemento che appartiene a $V$, che soddisfano determinate proprietà
        \begin{itemize}
            \item \textbf{somma vettoriale} di due vettori $\mathbf{u}$, $\mathbf{v}$ $\in V$: 
                \begin{equation}
                    \mathbf{u} + \mathbf{v} = \mathbf{w} \in V
                \end{equation}
            \item \textbf{moltiplicazione per uno scalare} di un vettore $\mathbf{u} \in V$ e uno scalare $a \in F$:
                \begin{equation}
                    a \mathbf{v} = \mathbf{w} \in V
                \end{equation}
        \end{itemize}
    \end{itemize}
\end{definition}

\begin{property}[Proprietà delle operazioni]
\end{property}

\begin{definition}[Base e dimensione di uno spazio]
\end{definition}


\subsection{Spazi vettoriali con prodotto interno}


\section{Applicazioni}
\subsection{Geometria}
\subsection{Fisica}


\chapter{Coordinate in spazi euclidei e cenni di calcolo vettoriale}
\begin{definition}[Vettore posizione]
\end{definition}

\begin{definition}[Coordinate]
\end{definition}
\begin{example}[Coordinate cartesiane]
\end{example}
\begin{example}[Coordinate polari]
\end{example}
\begin{example}[Coordinate sferiche e superficie terrestre]
\end{example}



\part{Geometria}

Introduzione storica:
\begin{itemize}
  \item Euclide
  \item Cartesio
  \item Riemann
\end{itemize}

\chapter{Geometria nel piano}
\section{Geometria euclidea}
\subsection{Introduzione}
\subsection{Rette e angoli}
\subsection{Triangoli}
\subsection{Circonferenza}

\section{Geometria cartesiana}
\subsection{Coordinate cartesiane}
\subsection{Punto, distanze, retta}
\begin{definition}[Punto]
Dato un sistema di coordinate cartesiane, un punto $P$ nel piano è individuato dalle sue due coordinate $(x,y)$.
\end{definition}
\begin{definition}[Distanza tra due punti]
La distanza tra due punti nel piano viene calcolata usando il teorema di Pitagora
\begin{equation}
    d_{PQ}^2 = (x_Q - x_P)^2 + (y_Q - y_P)^2 \ .
\end{equation}
\end{definition}
{\color{red} Perchè la distanza è data come una definizione? Geometria di Riemann: la distanza definisce tutte le proprietà di una geometria}

\begin{definition}[Retta]
    La retta può essere definita come l'insieme di punti $(x,y)$ equidistanti da due punti dati $P_1(x_1,y_1)$, $P_2(x_2,y_2)$.
\end{definition}
Partendo dalla definizione
\begin{equation}
\begin{aligned}
    d_1^2 & = d_2^2 \\
    (x - x_1)^2 + (y-y_1)^2 & = (x - x_2)^2 + (y-y_2)^2 \\
    x^2 - 2 x x_1 + x_1^2 + y^2 - 2 y y_1 + y_1^2 & = x^2 - 2 x x_2 + x_2^2 + y^2 - 2 y y_2 + y_2^2 \\
\end{aligned}
\end{equation}
\begin{equation}
  \qquad \rightarrow \qquad 2(x_2 - x_1) x + 2(y_2 - y_1) y + x_1^2 + y_1^2 - x_2^2 - y_2^2 = 0 \ .
\end{equation}
Quindi l'equazione generale della retta può essere riscritta nella forma
\begin{equation}
    A x + B y + C = 0 \ .
\end{equation}

\subsection{Trasformazioni di coordinate cartesiane e trasformazioni di curve}
\subsubsection{Traslazione dell'origine}
Sistema di coordinate $O'x'y'$ con assi paralleli al sistema di coordinate $Oxy$ e coordinate dell'origine $x_{O'}$, $y_{O'}$
\begin{equation}
    \begin{cases}
        x' = x - x_{O'} \\
        y' = y - y_{O'} \\
    \end{cases}
    \qquad \qquad
    \begin{cases}
        x  = x'+ x_{O'} \\
        y  = y'+ y_{O'} \\
    \end{cases}
\end{equation}
\subsubsection{Rotazione attorno all'origine}
Sistema di coordinate $O'x'y'$ origine coincidente con quella del sistema di coordinate $Oxy$ e assi rotati di un angolo $\theta$
\begin{equation}
    \begin{cases}
        x' =  x \cos \theta + y \sin \theta \\
        y' = -x \sin \theta + y \cos \theta \\
    \end{cases}
    \qquad \qquad
    \begin{cases}
        x  =  x'\cos \theta - y'\sin \theta \\
        y  =  x'\sin \theta + y'\cos \theta \\
    \end{cases}
\end{equation}

\subsection{Coniche}
\subsubsection{Parabola}
\begin{definition}[Parabola] Insieme dei punti $P$ del piano equidistanti da un punto $F$, chiamato \textbf{fuoco}, e una retta $r$ chiamata \textbf{direttrice}.
\begin{equation}
    \text{dist}(P,F) = \text{dist}(P,r) 
\end{equation}
\end{definition}
Scegliendo il fuoco $F(0, d)$ e la retta $r: y=-d$, si ricava l'equazione della parabola con vertice nell'origine e asse coincidente con l'asse $y$ degli assi cartesiani.
\begin{equation}
\begin{aligned}
    d^2_{PF} & = d^2_{Pr} \\
    (x - x_F)^2 + (y - y_F)^2 & = (y-y_r)^2 \\
    x^2 + (y - d)^2 & = (y+d)^2 \\
    x^2 + y^2 - 2dy + d^2 & = y^2 + 2dy + d^2 
\end{aligned}
\end{equation}
\begin{equation}
    \rightarrow \qquad 4d y = x^2 \qquad \rightarrow \qquad y = \dfrac{1}{4d} x^2 \ .
\end{equation}

\subsubsection{Ellisse}
\begin{definition}[Ellisse] Insieme dei punti $P$ del piano la cui somma delle distanze da due punti $F_1$, $F_2$, chiamati \textbf{fuochi} dell'ellisse, è costante.
    \begin{equation}
        \text{dist}(P, F_1) + \text{dist}(P, F_2) = 2a
    \end{equation}
\end{definition}
Scegliendo i fuochi $F_1(-c,0)$, $F_2(c,0)$
\begin{equation}
\begin{aligned}
    & \sqrt{( x - x_1 )^2 + (y-y_1)^2} + \sqrt{(x-x_2)^2 + (y-y_2)^2} = 2a \\
    & \sqrt{( x + c )^2 + y^2} = 2a - \sqrt{(x-c)^2 + y^2} \\
    & x^2 + 2 c x + c^2 + y^2  = 4a^2 - 4 a \sqrt{(x-c)^2 + y^2} +
    x^2 - 2 c x + c^2 + y^2 \\
    & 4c x = 4a^2 - 4a  \sqrt{(x-c)^2 + y^2} \\
    & (cx - a^2)^2 = ( - a  \sqrt{(x-c)^2 + y^2})^2 \\
    & c^2x^2 - 2a^2cx + a^4 = a^2(x-c)^2 + a^2 y^2 \\
    & (a^2-c^2) x^2 + a^2y^2 = a^2 (a^2-c^2) 
\end{aligned}
\end{equation}
Definendo $b^2 := a^2 - c^2 > 0$, si può riscrivere l'equazione dell'ellisse con il centro nell'origine e gli assi coincidenti con gli assi cartesiani come
\begin{equation}
  \dfrac{x^2}{a^2} + \dfrac{y^2}{b^2} = 1 \ .
\end{equation}


\subsubsection{Iperbole}
\begin{definition}[Iperbole] Insieme dei punti $P$ del piano la cui differenza delle distanze da due punti $F_1$, $F_2$, chiamati \textbf{fuochi} dell'ellisse, è costante in valore assoluto.
    \begin{equation}
       | \text{dist}(P, F_1) - \text{dist}(P, F_2) | = 2a
    \end{equation}
\end{definition}
Scegliendo i fuochi $F_1(-c,0)$, $F_2(c,0)$
\begin{equation}
\begin{aligned}
    & \sqrt{( x - x_1 )^2 + (y-y_1)^2} - \sqrt{(x-x_2)^2 + (y-y_2)^2} = \mp 2a \\
    & \sqrt{( x + c )^2 + y^2} = \mp 2a + \sqrt{(x-c)^2 + y^2} \\
    & x^2 + 2 c x + c^2 + y^2  = 4a^2 \mp 4 a \sqrt{(x-c)^2 + y^2} +
    x^2 - 2 c x + c^2 + y^2 \\
    & 4c x = 4a^2 \mp 4a  \sqrt{(x-c)^2 + y^2} \\
    & (cx - a^2)^2 = ( \mp a  \sqrt{(x-c)^2 + y^2})^2 \\
    & c^2x^2 - 2a^2cx + a^4 = a^2(x-c)^2 + a^2 y^2 \\
    & (c^2-a^2) x^2 - a^2y^2 = a^2 (c^2-a^2) 
\end{aligned}
\end{equation}
Definendo $b^2 := c^2 - a^2 > 0$, si può riscrivere l'equazione dell'ellisse con il centro nell'origine e gli assi coincidenti con gli assi cartesiani come
\begin{equation}
  \dfrac{x^2}{a^2} - \dfrac{y^2}{b^2} = 1 \ .
\end{equation}


\chapter{Geometria nello spazio}
\section{Geometria euclidea}
\section{Geometria cartesiana}



\part{Statistica}


%----------------------------------------------------------------------------------------
%	SECTIONING EXAMPLES CHAPTER
%----------------------------------------------------------------------------------------

\chapterimage{orange2.jpg} % Chapter heading image
\chapterspaceabove{6.75cm} % Whitespace from the top of the page to the chapter title on chapter pages
\chapterspacebelow{7.25cm} % Amount of vertical whitespace from the top margin to the start of the text on chapter pages

%------------------------------------------------

\chapter{Sectioning Examples}\index{Sectioning}

\section{Section Title}\index{Sectioning!Sections}

Lorem ipsum dolor sit amet, consectetur adipiscing elit\footnote{Footnote example text\ldots Lorem ipsum dolor sit amet, consectetur adipiscing elit. Praesent porttitor arcu luctus, imperdiet urna iaculis, mattis eros. Pellentesque iaculis odio vel nisl ullamcorper, nec faucibus ipsum molestie.}. Praesent porttitor arcu luctus, imperdiet urna iaculis, mattis eros. Pellentesque iaculis odio vel nisl ullamcorper, nec faucibus ipsum molestie. Sed dictum nisl non aliquet porttitor. Etiam vulputate arcu dignissim, finibus sem et, viverra nisl. Aenean luctus congue massa, ut laoreet metus ornare in. Nunc fermentum nisi imperdiet lectus tincidunt vestibulum at ac elit. Nulla mattis nisl eu malesuada suscipit.

Aliquam arcu turpis, ultrices sed luctus ac, vehicula id metus. Morbi eu feugiat velit, et tempus augue. Proin ac mattis tortor. Donec tincidunt, ante rhoncus luctus semper, arcu lorem lobortis justo, nec convallis ante quam quis lectus. Aenean tincidunt sodales massa, et hendrerit tellus mattis ac. Sed non pretium nibh. Donec cursus maximus luctus. Vivamus lobortis eros et massa porta porttitor.

\subsection{Subsection Title}\index{Sectioning!Subsections}

Fusce varius orci ac magna dapibus porttitor. In tempor leo a neque bibendum sollicitudin. Nulla pretium fermentum nisi, eget sodales magna facilisis eu. Praesent aliquet nulla ut bibendum lacinia. Donec vel mauris vulputate, commodo ligula ut, egestas orci. Suspendisse commodo odio sed hendrerit lobortis. Donec finibus eros erat, vel ornare enim mattis et. Donec finibus dolor quis dolor tempus consequat. Mauris fringilla dui id libero egestas, ut mattis neque ornare. Ut condimentum urna pharetra ipsum consequat, eu interdum elit cursus. Vivamus scelerisque tortor et nunc ultricies, id tincidunt libero pharetra. Aliquam eu imperdiet leo. Morbi a massa volutpat velit condimentum convallis et facilisis dolor.

In hac habitasse platea dictumst. Curabitur mattis elit sit amet justo luctus vestibulum. In hac habitasse platea dictumst. Pellentesque lobortis justo enim, a condimentum massa tempor eu. Ut quis nulla a quam pretium eleifend nec eu nisl. Nam cursus porttitor eros, sed luctus ligula convallis quis. Nam convallis, ligula in auctor euismod, ligula mauris fringilla tellus, et egestas mauris odio eget diam. Praesent sodales in ipsum eu dictum. Mauris interdum porttitor fringilla. Proin tincidunt sodales leo at ornare. Donec tempus magna non mauris gravida luctus. Cras vitae arcu vitae mauris eleifend scelerisque. Nam sem sapien, vulputate nec felis eu, blandit convallis risus. Pellentesque sollicitudin venenatis tincidunt. In et ipsum libero.

\subsubsection{Subsubsection Title} \index{Sectioning!Subsubsections}

Maecenas consectetur metus at tellus finibus condimentum. Proin arcu lectus, ultrices non tincidunt et, tincidunt ut quam. Integer luctus posuere est, non maximus ante dignissim quis. Nunc a cursus erat. Curabitur suscipit nibh in tincidunt sagittis. Nam malesuada vestibulum quam id gravida. Proin ut dapibus velit. Vestibulum eget quam quis ipsum semper convallis. Duis consectetur nibh ac diam dignissim, id condimentum enim dictum. Nam aliquet ligula eu magna pellentesque, nec sagittis leo lobortis. Aenean tincidunt dignissim egestas. Morbi efficitur risus ante, id tincidunt odio pulvinar vitae.

\paragraph{Paragraph Title}\index{Sectioning!Paragraphs} Nullam mollis tellus lorem, sed congue ipsum euismod a. Donec pulvinar neque sed ligula ornare sodales. Nulla sagittis vel lectus nec laoreet. Nulla volutpat malesuada turpis at ultricies. Ut luctus velit odio, sagittis volutpat erat aliquet vel. Donec ac neque eget neque volutpat mollis. Vestibulum viverra ligula et sapien bibendum, vel vulputate ex euismod. Curabitur nec velit velit. Aliquam vulputate lorem elit, id tempus nisl finibus sit amet. Curabitur ex turpis, consequat at lectus id, imperdiet molestie augue. Curabitur eu eros molestie purus commodo hendrerit. Quisque auctor ipsum nec mauris malesuada, non fringilla nibh viverra. Quisque gravida, metus quis semper pulvinar, dolor nisl suscipit leo, vestibulum volutpat ante justo ultrices diam. Sed id facilisis turpis, et aliquet eros.

In malesuada ullamcorper urna, sed dapibus diam sollicitudin non. Donec elit odio, accumsan ac nisl a, tempor imperdiet eros. Donec porta tortor eu risus consequat, a pharetra tortor tristique. Morbi sit amet laoreet erat. Morbi et luctus diam, quis porta ipsum. Quisque libero dolor, suscipit id facilisis eget, sodales volutpat dolor. Nullam vulputate interdum aliquam. Mauris id convallis erat, ut vehicula neque. Sed auctor nibh et elit fringilla, nec ultricies dui sollicitudin. Vestibulum vestibulum luctus metus venenatis facilisis. Suspendisse iaculis augue at vehicula ornare. Sed vel eros ut velit fermentum porttitor sed sed massa. Fusce venenatis, metus a rutrum sagittis, enim ex maximus velit, id semper nisi velit eu purus.

%------------------------------------------------

\section*{Unnumbered Section}

\subsection*{Unnumbered Subsection}

\subsubsection*{Unnumbered Subsubsection}

%----------------------------------------------------------------------------------------
%	IN-TEXT ELEMENT EXAMPLES CHAPTER
%----------------------------------------------------------------------------------------

\chapter{In-text Element Examples}

\section{Referencing Publications}\index{Citation}

This statement requires citation \cite{Smith:2022jd}; this one is more specific \cite[162]{Smith:2021qr}.

%------------------------------------------------

\section{Link Examples}\index{Links}

This is a URL link: \href{https://www.latextemplates.com}{LaTeX Templates}. This is an email link: \href{mailto:example@example.com}{example@example.com}. This is a monospaced URL link: \url{https://www.LaTeXTemplates.com}.

%------------------------------------------------

\section{Lists}\index{Lists}

Lists are useful to present information in a concise and/or ordered way.

\subsection{Numbered List}\index{Lists!Numbered List}

\begin{enumerate}
	\item First numbered item
	\begin{enumerate}
		\item First indented numbered item
		\item Second indented numbered item
		\begin{enumerate}
			\item First second-level indented numbered item
		\end{enumerate}
	\end{enumerate}
	\item Second numbered item
	\item Third numbered item
\end{enumerate}

\subsection{Bullet Point List}\index{Lists!Bullet Points}

\begin{itemize}
	\item First bullet point item
	\begin{itemize}
		\item First indented bullet point item
		\item Second indented bullet point item
		\begin{itemize}
			\item First second-level indented bullet point item
		\end{itemize}
	\end{itemize}
	\item Second bullet point item
	\item Third bullet point item
\end{itemize}

\subsection{Descriptions and Definitions}\index{Lists!Descriptions and Definitions}

\begin{description}
	\item[Name] Description
	\item[Word] Definition
	\item[Comment] Elaboration
\end{description}

%------------------------------------------------

\section{International Support}

àáâäãåèéêëìíîïòóôöõøùúûüÿýñçčšž

\noindent ÀÁÂÄÃÅÈÉÊËÌÍÎÏÒÓÔÖÕØÙÚÛÜŸÝÑ

\noindent ßÇŒÆČŠŽ

%------------------------------------------------

\section{Ligatures}

fi fj fl ffl ffi Ty

%----------------------------------------------------------------------------------------
%	PART
%----------------------------------------------------------------------------------------

\part{Part Two Title}

%----------------------------------------------------------------------------------------
%	MATHEMATICS EXAMPLES CHAPTER
%----------------------------------------------------------------------------------------

\chapter{Mathematics}

\section{Theorems}\index{Theorems}

\subsection{Several equations}\index{Theorems!Several Equations}

This is a theorem consisting of several equations.

\begin{theorem}[Name of the theorem] % Specify a name/title in square brackets, or leave them out for no title
	In $E=\mathbb{R}^n$ all norms are equivalent. It has the properties:
	\begin{align}
		& \big| ||\mathbf{x}|| - ||\mathbf{y}|| \big|\leq || \mathbf{x}- \mathbf{y}||\\
		&  ||\sum_{i=1}^n\mathbf{x}_i||\leq \sum_{i=1}^n||\mathbf{x}_i||\quad\text{where $n$ is a finite integer}
	\end{align}
\end{theorem}

\subsection{Single Line}\index{Theorems!Single Line}

This is a theorem consisting of just one line.

\begin{theorem} % Specify a name/title in square brackets, or leave them out for no title
	A set $\mathcal{D}(G)$ in dense in $L^2(G)$, $|\cdot|_0$. 
\end{theorem}

%------------------------------------------------

\section{Definitions}\index{Definitions}

A definition can be mathematical or it could define a concept.

\begin{definition}[Definition name] % Specify a name/title in square brackets, or leave them out for no title
	Given a vector space $E$, a norm on $E$ is an application, denoted $||\cdot||$, $E$ in $\mathbb{R}^+=[0,+\infty[$ such that:
	\begin{align}
		& ||\mathbf{x}||=0\ \Rightarrow\ \mathbf{x}=\mathbf{0}\\
		& ||\lambda \mathbf{x}||=|\lambda|\cdot ||\mathbf{x}||\\
		& ||\mathbf{x}+\mathbf{y}||\leq ||\mathbf{x}||+||\mathbf{y}||
	\end{align}
\end{definition}

%------------------------------------------------

\section{Notations}\index{Notations}

\begin{notation} % Specify a name/title in square brackets, or leave them out for no title
	Given an open subset $G$ of $\mathbb{R}^n$, the set of functions $\varphi$ are:
	\begin{enumerate}
		\item Bounded support $G$;
		\item Infinitely differentiable;
	\end{enumerate}
	a vector space is denoted by $\mathcal{D}(G)$. 
\end{notation}

%------------------------------------------------

\section{Remarks}\index{Remarks}

This is an example of a remark.

\begin{remark}
	The concepts presented here are now in conventional employment in mathematics. Vector spaces are taken over the field $\mathbb{K}=\mathbb{R}$, however, established properties are easily extended to $\mathbb{K}=\mathbb{C}$.
\end{remark}

%------------------------------------------------

\section{Corollaries}\index{Corollaries}

\begin{corollary}[Corollary name] % Specify a name/title in square brackets, or leave them out for no title
	The concepts presented here are now in conventional employment in mathematics. Vector spaces are taken over the field $\mathbb{K}=\mathbb{R}$, however, established properties are easily extended to $\mathbb{K}=\mathbb{C}$.
\end{corollary}

%------------------------------------------------

\section{Propositions}\index{Propositions}

\subsection{Several equations}\index{Propositions!Several Equations}

\begin{proposition}[Proposition name] % Specify a name/title in square brackets, or leave them out for no title
	It has the properties:
	\begin{align}
		& \big| ||\mathbf{x}|| - ||\mathbf{y}|| \big|\leq || \mathbf{x}- \mathbf{y}||\\
		&  ||\sum_{i=1}^n\mathbf{x}_i||\leq \sum_{i=1}^n||\mathbf{x}_i||\quad\text{where $n$ is a finite integer}
	\end{align}
\end{proposition}

\subsection{Single Line}\index{Propositions!Single Line}

\begin{proposition} % Specify a name/title in square brackets, or leave them out for no title
	Let $f,g\in L^2(G)$; if $\forall \varphi\in\mathcal{D}(G)$, $(f,\varphi)_0=(g,\varphi)_0$ then $f = g$. 
\end{proposition}

%------------------------------------------------

\section{Examples}\index{Examples}

\subsection{Equation Example}\index{Examples!Equation}

\begin{example} % Specify a name/title in square brackets, or leave them out for no title
	Let $G=\{x\in\mathbb{R}^2:|x|<3\}$ and denoted by: $x^0=(1,1)$; consider the function:
	\begin{equation}
	f(x)=\left\{\begin{aligned} & \mathrm{e}^{|x|} & & \text{si $|x-x^0|\leq 1/2$}\\
	& 0 & & \text{si $|x-x^0|> 1/2$}\end{aligned}\right.
	\end{equation}
	The function $f$ has bounded support, we can take $A=\{x\in\mathbb{R}^2:|x-x^0|\leq 1/2+\epsilon\}$ for all $\epsilon\in\mathopen{]}0\,;5/2-\sqrt{2}\mathclose{[}$.
\end{example}

\subsection{Text Example}\index{Examples!Text}

\begin{example}[Example name] % Specify a name/title in square brackets, or leave them out for no title
	Aliquam arcu turpis, ultrices sed luctus ac, vehicula id metus. Morbi eu feugiat velit, et tempus augue. Proin ac mattis tortor. Donec tincidunt, ante rhoncus luctus semper, arcu lorem lobortis justo, nec convallis ante quam quis lectus. Aenean tincidunt sodales massa, et hendrerit tellus mattis ac. Sed non pretium nibh. Donec cursus maximus luctus. Vivamus lobortis eros et massa porta porttitor.
\end{example}

%------------------------------------------------

\section{Exercises}\index{Exercises}

\begin{exercise} % Specify a name/title in square brackets, or leave them out for no title
	This is a good place to ask a question to test learning progress or further cement ideas into students' minds.
\end{exercise}

%------------------------------------------------

\section{Problems}\index{Problems}

\begin{problem} % Specify a name/title in square brackets, or leave them out for no title
	What is the average airspeed velocity of an unladen swallow?
\end{problem}

%------------------------------------------------

\section{Vocabulary}\index{Vocabulary}

Define a word to improve a students' vocabulary.

\begin{vocabulary}[Word] % Specify a name/title in square brackets, or leave them out for no title
	Definition of word.
\end{vocabulary}

%----------------------------------------------------------------------------------------
%	PRESENTING INFORMATION/RESULTS EXAMPLES CHAPTER
%----------------------------------------------------------------------------------------

\chapterimage{orange3.jpg} % Chapter heading image
\chapterspaceabove{6.25cm} % Whitespace from the top of the page to the chapter title on chapter pages
\chapterspacebelow{7.5cm} % Amount of vertical whitespace from the top margin to the start of the text on chapter pages

%------------------------------------------------

\chapter{Presenting Information and Results with a Long Chapter Title}

\section{Table}\index{Table}

Lorem ipsum dolor sit amet, consectetur adipiscing elit. Praesent porttitor arcu luctus, imperdiet urna iaculis, mattis eros. Pellentesque iaculis odio vel nisl ullamcorper, nec faucibus ipsum molestie. Sed dictum nisl non aliquet porttitor. Etiam vulputate arcu dignissim, finibus sem et, viverra nisl. Aenean luctus congue massa, ut laoreet metus ornare in. Nunc fermentum nisi imperdiet lectus tincidunt vestibulum at ac elit. Nulla mattis nisl eu malesuada suscipit.

\begin{table}[H] % Use [H] to suppress floating and place the figure/table exactly where it is specified in the text
	\centering % Horizontally center the table on the page
	\begin{tabular}{L{0.15\textwidth} R{0.15\textwidth} R{0.15\textwidth}} % Specify column alignment with L{width}, C{width} and R{width} for fixed-width columns, or the default latex l, c and r for flexible-width columns
		\toprule
		\textbf{Treatments} & \textbf{Response 1} & \textbf{Response 2}\\
		\midrule
		Treatment 1 & 0.0003262 & 0.562 \\
		Treatment 2 & 0.0015681 & 0.910 \\
		Treatment 3 & 0.0009271 & 0.296 \\
		\bottomrule
	\end{tabular}
	\caption{Table caption.}
	\label{tab:example} % Unique label used for referencing the table in-text
\end{table}

Referencing \autoref{tab:example} in-text using its label.

\begin{table}[t] % Floating table, [t] tells LaTeX to place it at the top of the next available page
	\centering % Horizontally center the table on the page
	\begin{tabular}{L{0.15\textwidth} R{0.15\textwidth} R{0.15\textwidth}} % Specify column alignment with L{width}, C{width} and R{width} for fixed-width columns, or the default latex l, c and r for flexible-width columns
		\toprule
		\textbf{Treatments} & \textbf{Response 1} & \textbf{Response 2}\\
		\midrule
		Treatment 1 & 0.0003262 & 0.562 \\
		Treatment 2 & 0.0015681 & 0.910 \\
		Treatment 3 & 0.0009271 & 0.296 \\
		\bottomrule
	\end{tabular}
	\caption{Floating table.}
	\label{tab:floating} % Unique label used for referencing the table in-text
\end{table}

%------------------------------------------------

\section{Figure}\index{Figure}

Lorem ipsum dolor sit amet, consectetur adipiscing elit. Praesent porttitor arcu luctus, imperdiet urna iaculis, mattis eros. Pellentesque iaculis odio vel nisl ullamcorper, nec faucibus ipsum molestie. Sed dictum nisl non aliquet porttitor. Etiam vulputate arcu dignissim, finibus sem et, viverra nisl. Aenean luctus congue massa, ut laoreet metus ornare in. Nunc fermentum nisi imperdiet lectus tincidunt vestibulum at ac elit. Nulla mattis nisl eu malesuada suscipit.

\begin{figure}[H] % Use [H] to suppress floating and place the figure/table exactly where it is specified in the text
	\centering % Horizontally center the figure on the page
	\includegraphics[width=0.5\textwidth]{creodocs_logo.pdf} % Include the figure image
	\caption{Figure caption.}
	\label{fig:placeholder} % Unique label used for referencing the figure in-text
\end{figure}

Referencing \autoref{fig:placeholder} in-text using its label.

\begin{figure}[b] % Floating figure, [b] tells LaTeX to place it at the bottom of the next available page
	\centering % Horizontally center the figure on the page
	\includegraphics[width=\textwidth]{creodocs_logo.pdf} % Include the figure image
	\caption{Floating figure.}
	\label{fig:floating} % Unique label used for referencing the figure in-text
\end{figure}

%----------------------------------------------------------------------------------------

\stopcontents[part] % Manually stop the 'part' table of contents here so the previous Part page table of contents doesn't list the following chapters

%----------------------------------------------------------------------------------------
%	BIBLIOGRAPHY
%----------------------------------------------------------------------------------------

\chapterimage{} % Chapter heading image
\chapterspaceabove{2.5cm} % Whitespace from the top of the page to the chapter title on chapter pages
\chapterspacebelow{2cm} % Amount of vertical whitespace from the top margin to the start of the text on chapter pages

%------------------------------------------------

\chapter*{Bibliography}
\markboth{\sffamily\normalsize\bfseries Bibliography}{\sffamily\normalsize\bfseries Bibliography} % Set the page headers to display a Bibliography chapter name
\addcontentsline{toc}{chapter}{\textcolor{ocre}{Bibliography}} % Add a Bibliography heading to the table of contents

\section*{Articles}
\addcontentsline{toc}{section}{Articles} % Add the Articles subheading to the table of contents

\printbibliography[heading=bibempty, type=article] % Output article bibliography entries

\section*{Books}
\addcontentsline{toc}{section}{Books} % Add the Books subheading to the table of contents

\printbibliography[heading=bibempty, type=book] % Output book bibliography entries

%----------------------------------------------------------------------------------------
%	INDEX
%----------------------------------------------------------------------------------------

\cleardoublepage % Make sure the index starts on an odd (right side) page
\phantomsection
\addcontentsline{toc}{chapter}{\textcolor{ocre}{Index}} % Add an Index heading to the table of contents
\printindex % Output the index

%----------------------------------------------------------------------------------------
%	APPENDICES
%----------------------------------------------------------------------------------------

\chapterimage{orange2.jpg} % Chapter heading image
\chapterspaceabove{6.75cm} % Whitespace from the top of the page to the chapter title on chapter pages
\chapterspacebelow{7.25cm} % Amount of vertical whitespace from the top margin to the start of the text on chapter pages

\begin{appendices}

\renewcommand{\chaptername}{Appendix} % Change the chapter name to Appendix, i.e. "Appendix A: Title", instead of "Chapter A: Title" in the headers

%------------------------------------------------

\chapter{Appendix Chapter Title}

\section{Appendix Section Title}

Lorem ipsum dolor sit amet, consectetur adipiscing elit. Aliquam auctor mi risus, quis tempor libero hendrerit at. Duis hendrerit placerat quam et semper. Nam ultricies metus vehicula arcu viverra, vel ullamcorper justo elementum. Pellentesque vel mi ac lectus cursus posuere et nec ex. Fusce quis mauris egestas lacus commodo venenatis. Ut at arcu lectus. Donec et urna nunc. Morbi eu nisl cursus sapien eleifend tincidunt quis quis est. Donec ut orci ex. Praesent ligula enim, ullamcorper non lorem a, ultrices volutpat dolor. Nullam at imperdiet urna. Pellentesque nec velit eget est euismod pretium.

%------------------------------------------------

\chapter{Appendix Chapter Title}

\section{Appendix Section Title}

Lorem ipsum dolor sit amet, consectetur adipiscing elit. Aliquam auctor mi risus, quis tempor libero hendrerit at. Duis hendrerit placerat quam et semper. Nam ultricies metus vehicula arcu viverra, vel ullamcorper justo elementum. Pellentesque vel mi ac lectus cursus posuere et nec ex. Fusce quis mauris egestas lacus commodo venenatis. Ut at arcu lectus. Donec et urna nunc. Morbi eu nisl cursus sapien eleifend tincidunt quis quis est. Donec ut orci ex. Praesent ligula enim, ullamcorper non lorem a, ultrices volutpat dolor. Nullam at imperdiet urna. Pellentesque nec velit eget est euismod pretium.

%------------------------------------------------

\end{appendices}

%----------------------------------------------------------------------------------------

\end{document}
